\chapter{Riemannian Geometry}\label{riemannian-geometry}
\chapterstatus{complete}

\section{Introduction}\label{riemannian-geometry:section-introduction}

A Riemannian metric turns a manifold into a metric space, produces a canonical connection (Theorem~\ref{riemannian-geometry:theorem-levi-civita})
and, together with an orientation, defines the Hodge star, the codifferential and the Laplacian
(Definition~\ref{riemannian-geometry:definition-laplacian}). Harmonic forms are the kernel of that Laplacian, and the Hodge theorem
(Chapter~\ref{hodge-theorem}) says that on a compact manifold they represent de Rham cohomology. Along the way we develop the local and global theory of geodesics: covariant derivatives along curves,
the Gauss lemma, the fact that geodesics are locally the shortest curves, and the Hopf--Rinow theorem, all with complete proofs; parallel transport and holonomy, with the holonomy principle, and the
theorems of Ambrose--Singer, de Rham and Berger; Jacobi fields, conjugate points and the first comparison theorems; the geometry of submanifolds, with the
first and second fundamental forms and the equations of Gauss and Codazzi; and the classical global theorems on surfaces: Gauss--Bonnet for regions with boundary,
and the theorems of Minding, Hadamard, Liebmann and Hilbert. References:
\cite{do-carmo-riemannian}, \cite{lee-riemannian}, \cite{kobayashi-nomizu}, \cite{joyce-special-holonomy}, \cite{warner-foundations}, \cite[Chapter 0]{griffiths-harris}.

\noindent\textbf{Prerequisites.} Chapter~\ref{connections} (Connections and Curvature).

\noindent\textbf{Conventions.} $(M, g)$ is a Riemannian manifold of dimension $n$, and we write $\langle v, w\rangle = g(v, w)$ and $|v| = \langle v, v\rangle^{1/2}$. Curves are piecewise smooth unless stated otherwise. In local
coordinates we use the summation convention: an index repeated once up and once down is summed over.

\section{Metrics and the Levi-Civita connection}\label{riemannian-geometry:section-levi-civita}

\begin{definition}\label{riemannian-geometry:definition-metric}
A \emph{Riemannian metric} on a smooth manifold $M$ is a smooth section $g$ of $\mathrm{Sym}^2T^*M$ that is an inner product on each tangent space; $(M, g)$ is a \emph{Riemannian manifold}. The
\emph{length} of a piecewise smooth curve is $L(\gamma) = \int_a^b|\dot\gamma(t)|\,dt$ with $|v| = g(v, v)^{1/2}$, and $d(p, q)$ is the infimum of the lengths of curves from $p$ to $q$. An \emph{isometry} is a
diffeomorphism preserving $g$.
\end{definition}

\begin{proposition}\label{riemannian-geometry:proposition-metric-exists}
Every smooth manifold admits a Riemannian metric, and $d$ is a metric inducing the topology of $M$.
\end{proposition}

\begin{proof}
Patch the Euclidean metrics of charts with a partition of unity (Theorem~\ref{smooth-manifolds:theorem-partitions-of-unity}); a convex combination of inner products is an inner product. Symmetry and the
triangle inequality for $d$ are clear, and $d(p, q) > 0$ for $p \neq q$ because in a chart around $p$ the metric is comparable to the Euclidean one on a compact neighbourhood, so curves leaving a small ball
have length bounded below. The same comparison shows that the $d$-balls form a basis of the topology.
\end{proof}

\begin{example}\label{riemannian-geometry:example-metrics}
The Euclidean metric on $\RR^n$; the induced metric on a submanifold, for instance the round metric on $S^n \subseteq \RR^{n+1}$; the hyperbolic metric $(dx^2 + dy^2)/y^2$ on the upper half plane; a
left-invariant metric on a Lie group, obtained from any inner product on $\mathfrak{g}$, which is bi-invariant when $G$ is compact (average over Haar measure,
Theorem~\ref{lie-groups:theorem-haar}); and the Fubini--Study metric on $\CC\PP^n$ (Chapter~\ref{kahler}).
\end{example}

\begin{proposition}\label{riemannian-geometry:proposition-quotient-metric}
Let a discrete group $G$ act freely and properly on $M$ by isometries of $g$. Then $M/G$ carries a unique Riemannian metric for which $\pi \colon M \to M/G$ is a local isometry.
\end{proposition}

\begin{proof}
By Theorem~\ref{smooth-manifolds:theorem-smooth-quotient}, $\pi$ is a local diffeomorphism. For $x \in M$ define the metric on $T_{\pi(x)}(M/G)$ by transporting $g_x$ through $d_x\pi$. This does not depend on the choice of $x$
in its orbit: $\pi \circ h = \pi$ for $h \in G$, so $d_{hx}\pi \circ d_xh = d_x\pi$, and $d_xh$ is an isometry $T_xM \to T_{hx}M$. On an open set $\pi(U)$ over which $\pi$ has the smooth inverse $(\pi|_U)^{-1}$ the metric is the pullback
of $g$ by this inverse, hence smooth. Uniqueness is clear.
\end{proof}

\begin{example}\label{riemannian-geometry:example-quotient-metrics}
The \emph{flat torus} $\RR^n/\ZZ^n$ and the \emph{flat Klein bottle} $\RR^2/G$ (Example~\ref{topological-manifolds:example-covering-actions}), since translations and the maps $(x, y) \mapsto (x + \frac12, -y)$ are Euclidean isometries; and
the round metric on $\RR\PP^n = S^n/\{\pm1\}$, as $-\mathrm{id}$ is an isometry of $S^n$.
\end{example}

\begin{theorem}[Fundamental theorem of Riemannian geometry]\label{riemannian-geometry:theorem-levi-civita}
On a Riemannian manifold there is a unique connection $\nabla$ on $TM$ that is metric ($Xg(Y, Z) = g(\nabla_XY, Z) + g(Y, \nabla_XZ)$) and torsion-free ($\nabla_XY - \nabla_YX = [X, Y]$): the
\emph{Levi-Civita connection}. It is determined by the Koszul formula
\[
2g(\nabla_XY, Z) = Xg(Y, Z) + Yg(X, Z) - Zg(X, Y) + g([X, Y], Z) - g([X, Z], Y) - g([Y, Z], X).
\]
In coordinates, $\nabla_{\partial_i}\partial_j = \sum_k\Gamma^k_{ij}\partial_k$ with $\Gamma^k_{ij} = \frac12\sum_lg^{kl}(\partial_ig_{jl} + \partial_jg_{il} - \partial_lg_{ij})$.
\end{theorem}

\begin{proof}
If $\nabla$ is metric and torsion-free, expanding the three terms $Xg(Y, Z)$, $Yg(X, Z)$, $-Zg(X, Y)$ with the metric property and using torsion-freeness to replace $\nabla_XY - \nabla_YX$ by $[X, Y]$ gives the
Koszul formula; since $g$ is nondegenerate, the formula determines $\nabla_XY$, which proves uniqueness. Conversely, define $\nabla_XY$ by the Koszul formula. The right-hand side is $\RR$-bilinear,
$C^\infty$-linear in $X$ and $Z$ and satisfies the Leibniz rule in $Y$ (the terms with derivatives of $f$ in $Y \mapsto fY$ combine to $2(Xf)g(Y, Z)$), so $\nabla$ is a connection; exchanging $X$ and $Y$ shows it
is torsion-free, and adding the formulas for $g(\nabla_XY, Z)$ and $g(Y, \nabla_XZ)$ gives $Xg(Y, Z)$. The coordinate formula follows by taking $X, Y, Z$ coordinate fields, whose brackets vanish.
\end{proof}

In coordinates, torsion-freeness says $\Gamma^k_{ij} = \Gamma^k_{ji}$, and the metric property for coordinate fields says
\[
\partial_ig_{jk} = \Gamma^l_{ij}g_{lk} + \Gamma^l_{ik}g_{jl}. \tag{$*$}
\]

\begin{definition}\label{riemannian-geometry:definition-geodesic}
A \emph{geodesic} is a curve with $\nabla_{\dot\gamma}\dot\gamma = 0$; in coordinates, $\ddot\gamma^k + \sum_{i,j}\Gamma^k_{ij}\dot\gamma^i\dot\gamma^j = 0$. By
Theorem~\ref{real-analysis:theorem-picard-lindelof} there is a unique maximal geodesic with given initial position and velocity, depending smoothly on them; the \emph{exponential map}
$\exp_p(v) = \gamma_v(1)$ is defined on a neighbourhood of $0 \in T_pM$.
\end{definition}

\section{Covariant derivatives along curves}\label{riemannian-geometry:section-along-curves}

A \emph{vector field along} a smooth curve $\gamma \colon I \to M$ is a smooth map $V \colon I \to TM$ with $V(t) \in T_{\gamma(t)}M$; for instance the velocity $\dot\gamma$, or $Y \circ \gamma$ for a vector field $Y$ on $M$.

\begin{proposition}\label{riemannian-geometry:proposition-covariant-derivative-curves}
There is a unique operator $D_t$ sending vector fields along $\gamma$ to vector fields along $\gamma$ such that $D_t$ is $\RR$-linear, $D_t(fV) = \dot fV + fD_tV$ for smooth $f \colon I \to \RR$, and
$D_t(Y \circ \gamma) = \nabla_{\dot\gamma}Y$ for every vector field $Y$ defined near $\gamma(I)$. In coordinates, for $V = V^k\partial_k$,
\[
D_tV = \big(\dot V^k + \Gamma^k_{ij}(\gamma)\dot\gamma^iV^j\big)\partial_k.
\]
Moreover $\frac{d}{dt}\langle V, W\rangle = \langle D_tV, W\rangle + \langle V, D_tW\rangle$. A vector field along $\gamma$ is parallel in the sense of Definition~\ref{connections:definition-parallel} if and only if $D_tV = 0$.
\end{proposition}

\begin{proof}
In a chart, $V = V^k(\partial_k \circ \gamma)$, and the axioms force $D_tV = \dot V^k\partial_k + V^k\nabla_{\dot\gamma}\partial_k = (\dot V^k + \dot\gamma^iV^j\Gamma^k_{ij})\partial_k$, since $\nabla_{\dot\gamma}\partial_j = \dot\gamma^i\Gamma^k_{ij}\partial_k$. This
proves uniqueness. Conversely, the formula satisfies the axioms in each chart (for $Y = Y^k\partial_k$, $\frac{d}{dt}Y^k(\gamma) = \dot\gamma^i\partial_iY^k$ matches the coordinate expression of $\nabla_{\dot\gamma}Y$), and by uniqueness the
local definitions agree on overlaps. For the metric property, differentiate $\langle V, W\rangle = g_{jk}(\gamma)V^jW^k$: the derivative of $g_{jk}(\gamma)$ is $\dot\gamma^i\partial_ig_{jk}$, and by $(*)$ the result equals
$\langle D_tV, W\rangle + \langle V, D_tW\rangle$. The last statement is the local expression of $\nabla_{\dot\gamma}s = 0$ in the coordinate frame.
\end{proof}

\begin{lemma}[Symmetry lemma]\label{riemannian-geometry:lemma-symmetry}
Let $f \colon (-\varepsilon, \varepsilon) \times (a, b) \to M$ be smooth, and write $D_s$, $D_t$ for the covariant derivatives along the curves $s \mapsto f(s, t)$ and $t \mapsto f(s, t)$. Then $D_s\partial_tf = D_t\partial_sf$.
\end{lemma}

\begin{proof}
In coordinates both sides equal $(\partial_s\partial_tf^k + \Gamma^k_{ij}\partial_sf^i\partial_tf^j)\partial_k$, by the symmetry $\Gamma^k_{ij} = \Gamma^k_{ji}$.
\end{proof}

\begin{proposition}\label{riemannian-geometry:proposition-geodesic-speed}
A curve $\gamma$ is a geodesic (Definition~\ref{riemannian-geometry:definition-geodesic}) if and only if $D_t\dot\gamma = 0$. Geodesics have constant speed, and parallel transport along any curve is a linear isometry.
\end{proposition}

\begin{proof}
The coordinate expression of $D_t\dot\gamma$ is the geodesic equation. If $D_t\dot\gamma = 0$ then $\frac{d}{dt}|\dot\gamma|^2 = 2\langle D_t\dot\gamma, \dot\gamma\rangle = 0$. If $V$ and $W$ are parallel, $\langle V, W\rangle$ is constant by
Proposition~\ref{riemannian-geometry:proposition-covariant-derivative-curves}.
\end{proof}

\section{The exponential map and minimising geodesics}\label{riemannian-geometry:section-exponential}

Geodesics are the projections to $M$ of the integral curves of a vector field on $TM$, the \emph{geodesic spray}, which in the coordinates $(x, v)$ of $TM$ is $v^k\partial_{x^k} - \Gamma^k_{ij}(x)v^iv^j\partial_{v^k}$. By
Theorem~\ref{smooth-manifolds:theorem-flow}, the set $\mathcal{E} \subseteq TM$ of vectors $v$ for which $\gamma_v$ is defined on $[0, 1]$ is open, and $\exp \colon \mathcal{E} \to M$ is smooth. Since $\gamma_{cv}(t) = \gamma_v(ct)$, each
$\mathcal{E}_p = \mathcal{E} \cap T_pM$ is star-shaped about $0$, and $\exp_p(tv) = \gamma_v(t)$.

\begin{proposition}\label{riemannian-geometry:proposition-exponential}
$d_0\exp_p = \id$, so $\exp_p$ is a diffeomorphism from a neighbourhood of $0$ onto a neighbourhood of $p$, giving \emph{normal coordinates} in which $g_{ij}(p) = \delta_{ij}$ and $\Gamma^k_{ij}(p) = 0$.
Geodesics through $p$ are the images of lines through $0$.
\end{proposition}

\begin{proof}
Since $\exp_p(tv) = \gamma_v(t)$, the differential of $\exp_p$ at $0$ sends $v$ to $\dot\gamma_v(0) = v$, and the inverse function theorem applies. In normal coordinates for an orthonormal basis of $T_pM$, $g_{ij}(p) = \delta_{ij}$, and the lines
$t \mapsto tv$ are geodesics, so the geodesic equation at $t = 0$ gives $\Gamma^k_{ij}(p)v^iv^j = 0$ for all $v$; as $\Gamma^k_{ij}$ is symmetric in $i$, $j$, it vanishes at $p$.
\end{proof}

\begin{theorem}[Gauss lemma]\label{riemannian-geometry:theorem-gauss-lemma}
Let $v \in \mathcal{E}_p$ and $w \in T_pM$, and identify $T_v(T_pM)$ with $T_pM$. Then $\langle d_v\exp_p(v), d_v\exp_p(w)\rangle = \langle v, w\rangle$.
\end{theorem}

\begin{proof}
Since $\mathcal{E}_p$ is open and star-shaped, $f(s, t) = \exp_p(t(v + sw))$ is defined and smooth for $|s|$ small and $t \in [0, 1]$. Let $T = \partial_tf$ and $S = \partial_sf$. For fixed $s$, $t \mapsto f(s, t)$ is the geodesic with initial
velocity $v + sw$, so $D_tT = 0$ and $|T|^2 = |v + sw|^2$ for all $t$. Using Proposition~\ref{riemannian-geometry:proposition-covariant-derivative-curves} and Lemma~\ref{riemannian-geometry:lemma-symmetry},
\[
\partial_t\langle S, T\rangle = \langle D_tS, T\rangle + \langle S, D_tT\rangle = \langle D_sT, T\rangle = \tfrac12\partial_s|T|^2 = \tfrac12\partial_s|v + sw|^2 = \langle v + sw, w\rangle.
\]
At $s = 0$ this is $\langle v, w\rangle$. Also $f(s, 0) = p$, so $S(s, 0) = 0$. Integrating, $\langle S, T\rangle(0, 1) = \langle v, w\rangle$. Finally $T(0, 1) = \frac{d}{dt}\exp_p(tv)|_{t=1} = d_v\exp_p(v)$ and
$S(0, 1) = \frac{d}{ds}\exp_p(v + sw)|_{s=0} = d_v\exp_p(w)$.
\end{proof}

\begin{theorem}[Geodesics are locally minimising]\label{riemannian-geometry:theorem-local-minimisation}
Let $\varepsilon > 0$ be such that $\exp_p$ is a diffeomorphism from $B(0, \varepsilon) \subseteq T_pM$ onto an open set $U$ (a \emph{normal ball}).
\begin{enumerate}
\item For $|v| < \varepsilon$, every curve from $p$ to $\exp_p(v)$ has length $\geq |v|$, with equality if and only if it is a monotone reparametrisation of the radial geodesic $t \mapsto \exp_p(tv)$, $t \in [0, 1]$.
\item $d(p, \exp_p(v)) = |v|$ for $|v| < \varepsilon$, and $\exp_p(B(0, r))$ is the metric ball $\{x : d(p, x) < r\}$ for $0 < r \leq \varepsilon$.
\end{enumerate}
\end{theorem}

\begin{proof}
For $x \in U$ let $\rho(x) = |\exp_p^{-1}(x)|$, a continuous function, smooth on $U \setminus \{p\}$. Let $c \colon [a, b] \to U \setminus \{p\}$ be a piecewise smooth curve and write $c(t) = \exp_p(\rho(t)u(t))$ with $\rho(t) = \rho(c(t)) > 0$ and $|u(t)| = 1$.
Then $\dot c = d\exp_p(\dot\rho u + \rho\dot u)$ at the point $\rho u$, with $\langle u, \dot u\rangle = 0$. By the Gauss lemma, $d\exp_p(u)$ has length $1$ and is orthogonal to $d\exp_p(\rho\dot u)$ (both at $\rho u$), so
\[
|\dot c|^2 = \dot\rho^2 + |d\exp_p(\rho\dot u)|^2 \geq \dot\rho^2, \qquad\text{hence}\qquad L(c) \geq \int_a^b|\dot\rho|\,dt \geq \rho(b) - \rho(a). \tag{$**$}
\]
Now let $c \colon [a, b] \to M$ go from $p$ to $q = \exp_p(v)$, $0 < |v| = r < \varepsilon$. Suppose first that $c$ stays in $U$. Let $a'$ be the last time with $c(a') = p$; applying $(**)$ on $[a'', b]$ for $a'' \downarrow a'$ gives $L(c) \geq r$. If
$L(c) = r$, then $c$ is constant on $[a, a']$, and on $(a', b]$ equality holds in $(**)$: $\dot\rho \geq 0$ and $d\exp_p(\rho\dot u) = 0$, so $\dot u = 0$ ($d\exp_p$ being injective on $B(0, \varepsilon)$), $u$ is constant and $c$ is a monotone
reparametrisation of the radial geodesic. Conversely such reparametrisations have length $r$. If $c$ leaves $U$, choose $r < r' < \varepsilon$. The set $K = \exp_p(\bar B(0, r'))$ is compact, hence closed, and contains $p$; so $c$
leaves $K$, and by continuity $\rho(c(t))$ takes the value $r'$ at some first time $t_2$ with $c([a, t_2]) \subseteq K$. The first case applied to $c|_{[a, t_2]}$ gives $L(c) \geq r' > r$. This proves (1).

For (2), the radial geodesic has length $|v|$, so $d(p, \exp_p(v)) = |v|$ by (1). The ball $\exp_p(B(0, r))$ lies in the metric ball of radius $r$. Conversely, if $d(p, x) < r$, choose a curve $c$ from $p$ to $x$ with $L(c) < r' < r$; if
$x \notin \exp_p(B(0, r'))$, the argument above shows $L(c) \geq r'$, a contradiction.
\end{proof}

\begin{remark}\label{riemannian-geometry:remark-convex}
Every point also has a \emph{geodesically convex} neighbourhood, in which any two points are joined by a unique minimising geodesic, contained in the neighbourhood \cite[Chapter 3]{do-carmo-riemannian}.
\end{remark}

\begin{lemma}[Uniformly normal neighbourhoods]\label{riemannian-geometry:lemma-uniformly-normal}
Every $q \in M$ has a neighbourhood $W$ and a $\delta > 0$ such that for every $x \in W$, $\exp_x$ is defined on $B(0, \delta) \subseteq T_xM$, is a diffeomorphism from it onto an open set, and $\exp_x(B(0, \delta)) \supseteq W$.
\end{lemma}

\begin{proof}
Consider $E \colon \mathcal{E} \to M \times M$, $E(v) = (\pi(v), \exp(v))$. In coordinates $(x, v)$ near $0_q \in TM$, its differential at $0_q$ is $\begin{pmatrix}I & 0\\ I & I\end{pmatrix}$, since $\exp(0_x) = x$ and $d_0\exp_x = \mathrm{id}$; so by the
inverse function theorem $E$ is a diffeomorphism from a neighbourhood $O$ of $0_q$ onto a neighbourhood of $(q, q)$. Shrinking, we may take $O = \{v \in T_xM : x \in W', |v| < \delta\}$ for a neighbourhood $W'$ of $q$ and some $\delta > 0$
(the norm is continuous on $TM$), and choose a neighbourhood $W \subseteq W'$ of $q$ with $W \times W \subseteq E(O)$. For $x \in W$, $\exp_x$ is injective on $B(0, \delta)$ with invertible differential, hence a diffeomorphism onto an open
set; and for $y \in W$, $(x, y) = E(v)$ with $v \in O$, that is $y = \exp_x(v)$ with $|v| < \delta$.
\end{proof}

\begin{lemma}\label{riemannian-geometry:lemma-minimising-geodesic}
Let $c \colon [a, b] \to M$ be a piecewise smooth curve of constant speed with $L(c) = d(c(a), c(b))$. Then $c$ is a geodesic; in particular it is smooth.
\end{lemma}

\begin{proof}
Every subarc of $c$ is also minimising, since otherwise replacing it by a shorter curve would shorten $c$. Let $t_0 \in [a, b]$, and take $W$ and $\delta$ for $q = c(t_0)$ as in Lemma~\ref{riemannian-geometry:lemma-uniformly-normal}. For
$t_1 < t_0 < t_2$ close to $t_0$ (or $t_1 = t_0$ at the endpoints), $c([t_1, t_2]) \subseteq W$ and $L(c|_{[t_1, t_2]}) < \delta$. Then $c|_{[t_1, t_2]}$ is a minimising curve from $x = c(t_1)$ to a point of the normal ball $\exp_x(B(0, \delta))$, so by
Theorem~\ref{riemannian-geometry:theorem-local-minimisation}(1) it is a monotone reparametrisation of a radial geodesic from $x$, and having constant speed, it is that geodesic, affinely reparametrised. So $c$ is a geodesic near every $t_0$.
\end{proof}

\begin{theorem}[Hopf--Rinow]\label{riemannian-geometry:theorem-hopf-rinow}
For a connected Riemannian manifold $(M, g)$ the following are equivalent: $(M, d)$ is a complete metric space; the geodesics are defined for all time; closed bounded subsets are compact. In that case any
two points are joined by a minimising geodesic. In particular compact Riemannian manifolds are geodesically complete.
\end{theorem}

\begin{proof}
We show more: the three conditions are also equivalent to (d) \emph{$\exp_p$ is defined on all of $T_pM$ for some $p$}; and (d) implies that every $q$ is joined to $p$ by a minimising geodesic.

\emph{Step 1: (d) gives minimising geodesics from $p$.} Let $q \neq p$, $r = d(p, q)$, and choose $0 < \delta < r$ with $\exp_p$ a diffeomorphism on $B(0, \delta')$ for some $\delta' > \delta$. The \emph{geodesic sphere}
$S = \exp_p(\{|w| = \delta\})$ is compact; let $x_0 = \exp_p(\delta u) \in S$, $|u| = 1$, minimise $d(\cdot, q)$ on $S$. Every curve from $p$ to $q$ meets $S$ (by the argument in the proof of
Theorem~\ref{riemannian-geometry:theorem-local-minimisation}), and points of $S$ have distance $\delta$ from $p$, so $r = \delta + d(x_0, q)$. Let $\gamma(t) = \exp_p(tu)$, defined for all $t$ by (d), and let
$T = \{t \in [\delta, r] : d(\gamma(t), q) = r - t\}$, a closed set containing $\delta$. Let $t_0 = \max T$ and suppose $t_0 < r$. Apply the same construction at $\gamma(t_0)$: for small $\eta > 0$ there is $x_1$ on the geodesic sphere of
radius $\eta$ about $\gamma(t_0)$ with $d(\gamma(t_0), q) = \eta + d(x_1, q)$, so $d(x_1, q) = r - t_0 - \eta$ and $d(p, x_1) \geq d(p, q) - d(x_1, q) = t_0 + \eta$. The curve going along $\gamma$ from $p$ to $\gamma(t_0)$ and then along the
radial geodesic to $x_1$ has length $t_0 + \eta$, so it is minimising, and by Lemma~\ref{riemannian-geometry:lemma-minimising-geodesic} (after parametrising by arc length) it is a geodesic, hence equal to $\gamma$ on $[0, t_0 + \eta]$.
So $x_1 = \gamma(t_0 + \eta)$ and $t_0 + \eta \in T$, a contradiction. Hence $t_0 = r$, $\gamma(r) = q$ and $L(\gamma|_{[0, r]}) = r$.

\emph{Step 2: the equivalences.} If closed bounded sets are compact, a Cauchy sequence, being bounded, has a convergent subsequence, hence converges; so $M$ is complete. If $M$ is complete, let $\gamma \colon [0, b) \to M$ be a maximal unit-speed
geodesic with $b < \infty$. For $t_k \uparrow b$, $d(\gamma(t_k), \gamma(t_l)) \leq |t_k - t_l|$, so $(\gamma(t_k))$ is Cauchy and converges to some $q$. Take $W$, $\delta$ for $q$ as in Lemma~\ref{riemannian-geometry:lemma-uniformly-normal}; then
geodesics with initial point in $W$ and speed $1$ are defined on $(-\delta, \delta)$. Choose $t_k > b - \delta$ with $\gamma(t_k) \in W$; the geodesic with initial velocity $\dot\gamma(t_k)$ is defined up to time $t_k + \delta > b$ and extends $\gamma$ by
uniqueness, a contradiction. So geodesics are defined for all time, which implies (d). Finally, assume (d) and let $K$ be closed and bounded, say $K \subseteq \{x : d(p, x) \leq R\}$. By Step 1, this ball is contained in
$\exp_p(\bar B(0, R))$, which is compact; so $K$ is compact. The last statements follow from Step 1 applied at every point, and from the compactness of closed bounded subsets of a compact space.
\end{proof}

\begin{counterexample}\label{riemannian-geometry:counterexample-incomplete}
Completeness is needed. In $M = \RR^2 \setminus \{0\}$ with the Euclidean metric, the geodesic $t \mapsto (1 - t, 0)$ is not defined at $t = 1$, the Cauchy sequence $(1/k, 0)$ has no limit, and the set $\{x \in M : |x| \leq 1\}$ is closed in $M$
and bounded but not compact. The points $(1, 0)$ and $(-1, 0)$ have distance $2$, the infimum of the lengths of curves avoiding the origin, but no
curve of length $2$ joins them in $M$, since in $\RR^2$ the only such curve is the segment through the origin.
\end{counterexample}

\begin{remark}\label{riemannian-geometry:remark-non-unique-geodesics}
Minimising geodesics need not be unique: on the round sphere every half great circle joins two antipodal points, and has length $\pi$. And geodesics stop minimising: a great circle arc longer than $\pi$ is a geodesic
but not the shortest curve between its endpoints. Both phenomena are explained by conjugate points (Section~\ref{riemannian-geometry:section-curvature}).
\end{remark}

\section{Parallel transport and holonomy}\label{riemannian-geometry:section-holonomy}

The \emph{holonomy group} $\mathrm{Hol}_p(g) \subseteq \GL(T_pM)$ of $(M, g)$ at $p$ is the holonomy group of the Levi-Civita connection (Definition~\ref{connections:definition-holonomy}), for piecewise smooth loops at $p$. The
\emph{restricted holonomy group} $\mathrm{Hol}^0_p(g)$ is the subgroup of transports along null-homotopic loops.

\begin{proposition}\label{riemannian-geometry:proposition-holonomy-basic}
Let $(M, g)$ be connected and $p, q \in M$.
\begin{enumerate}
\item $\mathrm{Hol}_p(g) \subseteq \mathrm{O}(T_pM, g_p)$.
\item For a curve $\sigma$ from $p$ to $q$, $\mathrm{Hol}_q(g) = P_\sigma\mathrm{Hol}_p(g)P_\sigma^{-1}$; so the holonomy groups at different points are conjugate.
\item $\mathrm{Hol}^0_p(g)$ is a normal subgroup of $\mathrm{Hol}_p(g)$, and $[\gamma] \mapsto P_\gamma^{-1}\mathrm{Hol}^0_p(g)$ is a surjective homomorphism $\pi_1(M, p) \to \mathrm{Hol}_p(g)/\mathrm{Hol}^0_p(g)$. In particular
$\mathrm{Hol}_p = \mathrm{Hol}^0_p$ if $M$ is simply connected.
\item $M$ is orientable if and only if $\mathrm{Hol}_p(g) \subseteq \mathrm{SO}(T_pM)$.
\item $\mathrm{Hol}^0_p(g)$ is trivial if and only if $g$ is flat, $R = 0$.
\end{enumerate}
\end{proposition}

\begin{proof}
(1) is Proposition~\ref{riemannian-geometry:proposition-geodesic-speed}. (2) Loops $\gamma$ at $p$ correspond bijectively to loops $\bar\sigma * \gamma * \sigma$ at $q$, with transport $P_\sigma P_\gamma P_\sigma^{-1}$
(Proposition~\ref{connections:proposition-parallel-transport}). (3) If $\beta$ is null-homotopic and $\alpha$ is a loop at $p$, then $\bar\alpha * \beta * \alpha$ is null-homotopic, with transport $P_\alpha P_\beta P_\alpha^{-1}$; so
$\mathrm{Hol}^0_p$ is normal. If $\gamma \simeq \gamma'$, then $\bar\gamma * \gamma'$ is null-homotopic and $P_{\gamma'}P_\gamma^{-1} \in \mathrm{Hol}^0_p$, so the map is well defined on homotopy classes, and it is multiplicative since
$P_{\gamma * \delta}^{-1} = P_\gamma^{-1}P_\delta^{-1}$. Surjectivity is clear. (4) If $M$ is oriented, a parallel frame along a curve is continuous, so the set of times at which it is positively oriented is open and closed; hence parallel
transport preserves orientations. Conversely, if $\mathrm{Hol}_p \subseteq \mathrm{SO}(T_pM)$, fix an orientation of $T_pM$ and transport it along curves from $p$; by (2) the result does not depend on the curve, and near each point the
transported frames along radial geodesics form a smooth local frame, so the orientation is continuous. (5) By Theorem~\ref{connections:theorem-flat-local-systems}, $R = 0$ if and only if parallel transport depends only on the
homotopy class of the path with fixed endpoints, which is equivalent to the triviality of transport around null-homotopic loops.
\end{proof}

\begin{theorem}[Holonomy principle]\label{riemannian-geometry:theorem-holonomy-principle}
Let $(M, g)$ be connected and let $E$ be a tensor bundle $T^{\otimes a} \otimes (T^*)^{\otimes b}$ of $M$ with the connection induced by $\nabla$. A tensor $\tau_p \in E_p$ is the value at $p$ of a parallel section of $E$ if and only if it is
invariant under $\mathrm{Hol}_p(g)$, acting on $E_p$ through its action on $T_pM$; the parallel section is then unique. Similarly, a subspace $V \subseteq T_pM$ is the fibre at $p$ of a parallel subbundle of $TM$ if and only if it is
$\mathrm{Hol}_p(g)$-invariant.
\end{theorem}

\begin{proof}
Parallel transport on $E$ is induced by that on $TM$. If $\tau$ is parallel, then along every loop $\gamma$ at $p$ the section $\tau \circ \gamma$ is parallel, so $P_\gamma\tau_p = \tau_p$. Conversely, if $\tau_p$ is invariant, define $\tau_q = P_\sigma\tau_p$
for any curve $\sigma$ from $p$ to $q$. If $\sigma'$ is another such curve, $P_{\sigma'}^{-1}P_\sigma$ is the transport along the loop $\sigma * \bar\sigma'$ at $p$, which lies in $\mathrm{Hol}_p$ and fixes $\tau_p$; so
$\tau_q$ is well defined. Near a point $q_0$, $\tau_q$ is the transport of $\tau_{q_0}$ along the radial geodesic from $q_0$ to $q$ in a normal ball, which depends smoothly on $q$ by smooth dependence of solutions of
differential equations; so $\tau$ is smooth. For every curve $c$, $\tau \circ c$ is the transport of $\tau_{c(0)}$ along $c$, hence parallel; so $\nabla\tau = 0$. Uniqueness holds because a parallel section is determined by its value at one
point of a connected manifold. The statement for subbundles is proved in the same way.
\end{proof}

\begin{example}\label{riemannian-geometry:example-holonomy}
\begin{enumerate}
\item On $\RR^n$ the Christoffel symbols vanish in the standard coordinates, so parallel transport is the identity and $\mathrm{Hol} = \{1\}$; the same holds for the flat torus, since parallel transport in the quotient is induced from
$\RR^n$ and translations have trivial differential.
\item For the flat Klein bottle $\RR^2/G$ (Example~\ref{riemannian-geometry:example-quotient-metrics}), a loop at $\pi(\tilde x)$ lifts to a curve from $\tilde x$ to $h\tilde x$ for some $h \in G$; transport in $\RR^2$ is the identity, and
$d\pi_{h\tilde x} = d\pi_{\tilde x} \circ (dh)^{-1}$, so the holonomy is $(dh)^{-1}$, the inverse of the linear part of $h$. Hence $\mathrm{Hol} = \{1, \mathrm{diag}(1, -1)\}$: the Klein bottle is flat, $\mathrm{Hol}^0 = \{1\}$, but
$\mathrm{Hol} \neq \{1\}$, and $\mathrm{Hol} \not\subseteq \mathrm{SO}(2)$, in accordance with non-orientability.
\item For the round sphere $S^2$, $\mathrm{Hol}_p = \mathrm{Hol}^0_p$ is a nontrivial subgroup of $\mathrm{SO}(2)$ by Proposition~\ref{riemannian-geometry:proposition-holonomy-basic}(3),(4),(5), as $S^2$ is simply connected, orientable and
not flat. Since $\mathrm{Hol}^0_p$ is a connected Lie subgroup (Theorem~\ref{riemannian-geometry:theorem-ambrose-singer} below), it is $\mathrm{SO}(2)$. Concretely, transport around the boundary of the octant
$\{x, y, z \geq 0\}$ rotates tangent vectors by $\pi/2$, the area enclosed.
\end{enumerate}
\end{example}

\begin{theorem}[Ambrose--Singer]\label{riemannian-geometry:theorem-ambrose-singer}
$\mathrm{Hol}_p(g)$ is a Lie subgroup of $\mathrm{O}(T_pM)$ with identity component $\mathrm{Hol}^0_p(g)$, and the Lie algebra of $\mathrm{Hol}_p(g)$ is spanned by the endomorphisms $P_\sigma^{-1} \circ R(P_\sigma u, P_\sigma v) \circ P_\sigma$ of
$T_pM$, for all curves $\sigma$ starting at $p$ and all $u, v \in T_pM$.
\end{theorem}

This is \cite[Chapter II]{kobayashi-nomizu}; see also \cite{joyce-special-holonomy}. The theorem makes Proposition~\ref{riemannian-geometry:proposition-holonomy-basic}(5) quantitative: curvature measures infinitesimal holonomy.

\begin{theorem}[de Rham decomposition]\label{riemannian-geometry:theorem-de-rham-decomposition}
If the holonomy representation of a connected Riemannian manifold on $T_pM$ is reducible, $T_pM = V_1 \oplus V_2$ with both summands $\mathrm{Hol}_p$-invariant and nonzero, then $M$ is locally isometric to a Riemannian product.
If moreover $M$ is complete and simply connected, it is globally isometric to a product $M_1 \times M_2$ of Riemannian manifolds with $T_pM_i = V_i$.
\end{theorem}

\begin{proof}[Sketch of proof]
By Theorem~\ref{riemannian-geometry:theorem-holonomy-principle}, $V_1$ and $V_2$ extend to parallel, mutually orthogonal distributions $D_1$ and $D_2$. A parallel distribution is involutive, as $[X, Y] = \nabla_XY - \nabla_YX$ lies in
$D_i$ when $X$, $Y$ do, so by the Frobenius theorem (Theorem~\ref{smooth-manifolds:theorem-frobenius}) there are local coordinates $(x, y)$ with $D_1$ spanned by the $\partial_x$ and $D_2$ by the $\partial_y$; the omitted verification is that in
these coordinates the metric has the product form $g_1(x) + g_2(y)$, which follows from the parallelism of $D_1$ and $D_2$. The global statement is \cite[Chapter IV]{kobayashi-nomizu}.
\end{proof}

\begin{theorem}[Berger]\label{riemannian-geometry:theorem-berger}
Let $(M, g)$ be a simply connected Riemannian manifold of dimension $n$ that is not locally a product (its holonomy representation is irreducible) and not locally symmetric ($\nabla R \neq 0$). Then $\mathrm{Hol}(g)$ is one
of: $\mathrm{SO}(n)$; $\mathrm{U}(m)$ with $n = 2m \geq 4$ (K\"ahler metrics); $\mathrm{SU}(m)$ with $n = 2m \geq 4$ (Ricci-flat K\"ahler, Calabi--Yau); $\mathrm{Sp}(m)$ with $n = 4m \geq 8$ (hyper-K\"ahler); $\mathrm{Sp}(m)\mathrm{Sp}(1)$ with
$n = 4m \geq 8$ (quaternionic K\"ahler); $G_2$ with $n = 7$; $\mathrm{Spin}(7)$ with $n = 8$.
\end{theorem}

This is \cite{berger-1955}; Berger's list also contained $\mathrm{Spin}(9)$ in dimension $16$, which was later shown to occur only for symmetric metrics. See \cite{joyce-special-holonomy} for the modern treatment and for the
existence of metrics with holonomy $G_2$ and $\mathrm{Spin}(7)$.

\begin{remark}\label{riemannian-geometry:remark-special-holonomy}
By the holonomy principle, the reduction of the holonomy group to a subgroup of $\mathrm{O}(n)$ is the same as the existence of a parallel tensor with that stabiliser. For instance $\mathrm{Hol} \subseteq \mathrm{U}(m)$ if and only
if there is a parallel orthogonal almost complex structure $J$, that is, $(M, g, J)$ is K\"ahler (Chapter~\ref{kahler}); and $\mathrm{Hol} \subseteq \mathrm{SU}(m)$ if and only if there is moreover a parallel complex volume form,
which for a simply connected K\"ahler manifold is equivalent to Ricci-flatness (Chapter~\ref{calabi-yau}). This is the dictionary between Riemannian holonomy and complex geometry that underlies Part~\ref{part-complex-geometry}.
\end{remark}

\section{Curvature}\label{riemannian-geometry:section-curvature}

\begin{definition}\label{riemannian-geometry:definition-curvature-tensor}
The \emph{Riemann curvature tensor} is the curvature of the Levi-Civita connection (Proposition~\ref{connections:proposition-curvature}),
$R(X, Y)Z = \nabla_X\nabla_YZ - \nabla_Y\nabla_XZ - \nabla_{[X,Y]}Z$, viewed as the $4$-tensor $R(X, Y, Z, W) = g(R(X, Y)Z, W)$. For a plane $\sigma \subseteq T_pM$ with orthonormal basis $u, v$, the
\emph{sectional curvature} is $K(\sigma) = R(u, v, v, u)$; the \emph{Ricci tensor} is $\mathrm{Ric}(X, Y) = \sum_iR(e_i, X, Y, e_i)$ for an orthonormal frame, and the \emph{scalar curvature} is
$s = \sum_i\mathrm{Ric}(e_i, e_i)$.
\end{definition}

\begin{proposition}\label{riemannian-geometry:proposition-curvature-symmetries}
$R$ is tensorial and satisfies $R(X, Y, Z, W) = -R(Y, X, Z, W) = -R(X, Y, W, Z) = R(Z, W, X, Y)$ and the first Bianchi identity
$R(X, Y)Z + R(Y, Z)X + R(Z, X)Y = 0$. The sectional curvatures determine $R$.
\end{proposition}

\begin{proof}
Tensoriality and the first antisymmetry are general facts about curvature. The second follows by polarising $R(X, Y, Z, Z) = 0$, which holds because the connection is metric:
$0 = XY g(Z,Z)/2 - \cdots$ expands to $g(R(X,Y)Z, Z)$. The first Bianchi identity follows from torsion-freeness and the Jacobi identity, and the symmetry in the pairs is a formal consequence of the other
three, obtained by adding the Bianchi identities for the four cyclic permutations. The last statement is the standard polarisation argument \cite[Chapter 4]{do-carmo-riemannian}.
\end{proof}

\begin{example}\label{riemannian-geometry:example-curvature}
Euclidean space has $R = 0$; the round sphere of radius $1$ has constant sectional curvature $1$ and $\mathrm{Ric} = (n-1)g$; hyperbolic space has constant curvature $-1$. A Lie group with a bi-invariant
metric has $R(X, Y)Z = \frac14[[X, Y], Z]$ for left-invariant fields, hence non-negative sectional curvature. A metric has constant sectional curvature $\kappa$ if and only if
$R(X, Y)Z = \kappa(\langle Y, Z\rangle X - \langle X, Z\rangle Y)$, since the right-hand side has the symmetries of Proposition~\ref{riemannian-geometry:proposition-curvature-symmetries} and sectional curvature $\kappa$.
\end{example}

\begin{lemma}\label{riemannian-geometry:lemma-curvature-families}
For $f$ as in Lemma~\ref{riemannian-geometry:lemma-symmetry} and a vector field $V$ along $f$, $D_sD_tV - D_tD_sV = R(\partial_sf, \partial_tf)V$.
\end{lemma}

\begin{proof}
Write $V = V^k(\partial_k \circ f)$. In $D_sD_tV$ the terms containing derivatives of the $V^k$ are symmetric in $s$ and $t$, so the left side is $V^k(D_sD_t - D_tD_s)(\partial_k \circ f)$. Now
$D_t(\partial_k \circ f) = \partial_tf^i(\nabla_{\partial_i}\partial_k) \circ f$, and
$D_sD_t(\partial_k \circ f) = \partial_s\partial_tf^i(\nabla_{\partial_i}\partial_k) \circ f + \partial_tf^i\partial_sf^j(\nabla_{\partial_j}\nabla_{\partial_i}\partial_k) \circ f$. Antisymmetrising and using $[\partial_i, \partial_j] = 0$ gives
$\partial_sf^j\partial_tf^iR(\partial_j, \partial_i)\partial_k = R(\partial_sf, \partial_tf)\partial_k$.
\end{proof}

\begin{definition}\label{riemannian-geometry:definition-jacobi}
A \emph{Jacobi field} along a geodesic $\gamma$ is a vector field $J$ along $\gamma$ with $D_t^2J + R(J, \dot\gamma)\dot\gamma = 0$. A point $\gamma(b)$ is \emph{conjugate} to $\gamma(a)$ along $\gamma$ if there is a nonzero Jacobi field
vanishing at $a$ and $b$.
\end{definition}

The Jacobi equation is a linear second order equation, so a Jacobi field is determined by $J(a)$ and $D_tJ(a)$, and the Jacobi fields along $\gamma$ form a vector space of dimension $2n$.

\begin{proposition}\label{riemannian-geometry:proposition-jacobi}
\begin{enumerate}
\item If $f(s, t)$ is a family of geodesics $t \mapsto f(s, t)$, then $J = \partial_sf(0, -)$ is a Jacobi field along $\gamma = f(0, -)$.
\item For $v \in \mathcal{E}_p$ and $w \in T_pM$, $J(t) = d_{tv}\exp_p(tw)$ is the Jacobi field along $\gamma(t) = \exp_p(tv)$ with $J(0) = 0$ and $D_tJ(0) = w$.
\item $d_v\exp_p$ is singular if and only if $\exp_p(v)$ is conjugate to $p$ along $t \mapsto \exp_p(tv)$, $t \in [0, 1]$.
\end{enumerate}
\end{proposition}

\begin{proof}
(1) With $T = \partial_tf$ and $S = \partial_sf$, Lemmas~\ref{riemannian-geometry:lemma-symmetry} and~\ref{riemannian-geometry:lemma-curvature-families} give
$D_tD_tS = D_tD_sT = D_sD_tT + R(T, S)T = R(T, S)T = -R(S, T)T$, since $D_tT = 0$. (2) Apply (1) to $f(s, t) = \exp_p(t(v + sw))$: $\partial_sf(0, t) = d_{tv}\exp_p(tw)$. It vanishes at $t = 0$, and
$D_tS(0, 0) = D_sT(0, 0) = \frac{d}{ds}(v + sw) = w$, the covariant derivative along the constant curve $s \mapsto p$ being the ordinary derivative in $T_pM$. (3) By (2) and uniqueness, the Jacobi fields vanishing at $0$ are exactly the
$d_{tv}\exp_p(tw)$, and at $t = 1$ this is $d_v\exp_p(w)$.
\end{proof}

\begin{example}\label{riemannian-geometry:example-jacobi-constant}
If $g$ has constant sectional curvature $\kappa$ and $\gamma$ has unit speed, then for a parallel unit field $E \perp \dot\gamma$, $R(E, \dot\gamma)\dot\gamma = \kappa E$ by Example~\ref{riemannian-geometry:example-curvature}, and the Jacobi fields
normal to $\gamma$ with $J(0) = 0$ are $J(t) = s_\kappa(t)E(t)$ with $s_\kappa(t) = t$, $\sin(\sqrt\kappa t)/\sqrt\kappa$ or $\sinh(\sqrt{-\kappa}t)/\sqrt{-\kappa}$ for $\kappa = 0$, $> 0$, $< 0$. On the unit sphere they vanish again at $t = \pi$: antipodal
points are conjugate, which explains Remark~\ref{riemannian-geometry:remark-non-unique-geodesics}. For $\kappa \leq 0$ there are no conjugate points.
\end{example}

\begin{theorem}[Cartan--Hadamard]\label{riemannian-geometry:theorem-cartan-hadamard}
Let $(M, g)$ be complete and connected with sectional curvature $K \leq 0$. Then for every $p$, $\exp_p \colon T_pM \to M$ is a covering map; if $M$ is simply connected, it is a diffeomorphism, so $M$ is diffeomorphic to $\RR^n$.
\end{theorem}

\begin{proof}[Sketch of proof]
Let $J$ be a Jacobi field along a geodesic $\gamma$ with $J(0) = 0$ and $D_tJ(0) \neq 0$, and $\phi = |J|^2$. Then
$\phi'' = 2|D_tJ|^2 - 2\langle R(J, \dot\gamma)\dot\gamma, J\rangle = 2|D_tJ|^2 - 2K(J, \dot\gamma)(|J|^2|\dot\gamma|^2 - \langle J, \dot\gamma\rangle^2) \geq 0$, where $K(J, \dot\gamma)$ is the sectional curvature of the plane they span (if they are
independent). So $\phi$ is convex with $\phi(0) = \phi'(0) = 0$ and $\phi''(0) > 0$, hence positive for $t > 0$: there are no conjugate points, and $\exp_p$ is a local diffeomorphism on $T_pM$ (defined everywhere by
Theorem~\ref{riemannian-geometry:theorem-hopf-rinow}) by Proposition~\ref{riemannian-geometry:proposition-jacobi}(3). The omitted verification is that a local diffeomorphism from a complete Riemannian manifold (here $T_pM$ with the
pulled back metric, whose radial lines are geodesics) that is a local isometry is a covering map \cite[Chapter 7]{do-carmo-riemannian}.
\end{proof}

\begin{theorem}[Bonnet--Myers]\label{riemannian-geometry:theorem-bonnet-myers}
Let $(M, g)$ be complete and connected with $\mathrm{Ric} \geq (n - 1)\kappa g$ for some $\kappa > 0$. Then $M$ has diameter at most $\pi/\sqrt\kappa$; hence $M$ is compact, and $\pi_1(M)$ is finite.
\end{theorem}

This is \cite[Chapter 9]{do-carmo-riemannian}. The proof computes the second variation of length along a minimising geodesic of length $L > \pi/\sqrt\kappa$ in the directions $\sin(\pi t/L)E_i(t)$, for a parallel orthonormal frame
$E_i \perp \dot\gamma$, and finds that their sum is negative, contradicting minimality. Finiteness of $\pi_1$ follows by applying the diameter bound to the universal cover with the pulled back metric, which is compact, so the
fibres of the covering are finite.

\begin{remark}\label{riemannian-geometry:remark-comparison}
Theorems~\ref{riemannian-geometry:theorem-cartan-hadamard} and~\ref{riemannian-geometry:theorem-bonnet-myers} are the simplest \emph{comparison theorems}, which compare a manifold with a space of constant curvature
through its Jacobi fields. The Rauch comparison theorem \cite[Chapter 10]{do-carmo-riemannian} and the Bishop--Gromov volume comparison \cite{lee-riemannian} are their refinements, and they lead to the sphere theorems and to
Gromov's compactness theorem.
\end{remark}

\begin{theorem}[Gauss--Bonnet]\label{riemannian-geometry:theorem-gauss-bonnet}
Let $M$ be a compact oriented Riemannian surface, with Gaussian curvature $K$ (the sectional curvature of its tangent planes) and area form $dA$. Then
\[
\int_MK\,dA = 2\pi\chi(M) .
\]
\end{theorem}

\begin{proof}
Let $J$ be the rotation by $+\pi/2$ in each tangent plane, for the metric and the orientation: $Je_1 = e_2$ for every oriented orthonormal frame $(e_1, e_2)$. It is a smooth endomorphism with $J^2 = -1$, so $TM$ becomes a
complex line bundle, with $i$ acting as $J$, whose underlying oriented real bundle is $TM$. Near a point choose an oriented orthonormal frame $(e_1, e_2)$ and define the $1$-form $\omega$ by $\nabla_Xe_1 = \omega(X)e_2$; since
$\nabla$ is metric, $\langle\nabla_Xe_1, e_1\rangle = 0$ and $\nabla_Xe_2 = -\omega(X)e_1$. Hence $\nabla J = J\nabla$, so $\nabla$ is a connection on the complex line bundle, and in the complex frame $e_1$ it reads
$\nabla e_1 = \omega \otimes Je_1 = i\omega \otimes e_1$: its connection form is $\theta = i\omega$ and its curvature is $F = d\theta = i\,d\omega$ (Proposition~\ref{connections:proposition-curvature}).

We compare $d\omega$ with $K$. From the definition of the curvature tensor,
\[
R(X, Y)e_1 = \nabla_X(\omega(Y)e_2) - \nabla_Y(\omega(X)e_2) - \omega([X, Y])e_2 = \big(X\omega(Y) - Y\omega(X) - \omega([X, Y])\big)e_2 = d\omega(X, Y)\,e_2,
\]
the terms $\mp\omega(X)\omega(Y)e_1$ cancelling. By Definition~\ref{riemannian-geometry:definition-curvature-tensor} and Proposition~\ref{riemannian-geometry:proposition-curvature-symmetries},
$K = R(e_1, e_2, e_2, e_1) = -R(e_1, e_2, e_1, e_2) = -\langle R(e_1, e_2)e_1, e_2\rangle = -d\omega(e_1, e_2)$. So $d\omega = -K\,e^1 \wedge e^2 = -K\,dA$ and $F = -iK\,dA$.

By Chern--Weil theory (Theorem~\ref{characteristic-classes:theorem-chern-weil} and Definition~\ref{characteristic-classes:definition-chern-forms}), $c_1(TM)$ is represented by $\frac{i}{2\pi}F = \frac{1}{2\pi}K\,dA$ for this connection.
For a complex line bundle the Euler class of the underlying oriented real bundle is $c_1$ (Theorem~\ref{characteristic-classes:theorem-euler}(2)), and $\int_Me(TM) = \chi(M)$ by
Theorem~\ref{characteristic-classes:theorem-euler}(3). Hence $\frac{1}{2\pi}\int_MK\,dA = \chi(M)$.
\end{proof}

\begin{corollary}\label{riemannian-geometry:corollary-gauss-bonnet}
Let $\Sigma_g$ be the closed orientable surface of genus $g$, so that $\chi(\Sigma_g) = 2 - 2g$ (Example~\ref{singular-homology:example-cellular-computations}).
\begin{enumerate}
\item The integral of the Gaussian curvature does not depend on the metric: for the round unit sphere, $K = 1$ and the area is $4\pi = 2\pi \cdot 2$.
\item If $\Sigma_g$ carries a metric with $K > 0$ everywhere, then $g = 0$; if it carries a metric with $K \geq 0$ everywhere, then $g \leq 1$; if it carries one with $K < 0$ everywhere, then $g \geq 2$.
\item The torus admits flat metrics, as $\RR^2/\ZZ^2$, but no metric with $K > 0$ everywhere or $K < 0$ everywhere.
\end{enumerate}
\end{corollary}

\begin{proof}
(1) is Theorem~\ref{riemannian-geometry:theorem-gauss-bonnet}. (2) The sign of $\int_MK\,dA = 2\pi(2 - 2g)$ is positive, non-negative or negative accordingly, and $2 - 2g > 0$ forces $g = 0$, $2 - 2g \geq 0$ forces
$g \leq 1$, and $2 - 2g < 0$ forces $g \geq 2$. (3) The torus has $\chi = 0$, so $K$ cannot have constant nonzero sign, while the flat metric induced from $\RR^2$ has $K = 0$ (Example~\ref{riemannian-geometry:example-curvature}).
\end{proof}

\section{Submanifolds and the second fundamental form}\label{riemannian-geometry:section-submanifolds}

Classical differential geometry studies curves and surfaces in $\RR^3$ through two quadratic forms: the \emph{first fundamental form}, the metric induced from
$\RR^3$, and the \emph{second fundamental form}, which measures how the surface bends in the ambient space. This section develops both for a submanifold of
any Riemannian manifold, proves the equations of Gauss and Codazzi that relate them to the curvature, and then specialises to surfaces in $\RR^3$, where Gauss's
equation becomes the Theorema Egregium.

Throughout, $(\overline{M}, \bar{g})$ is a Riemannian manifold with Levi-Civita connection $\overline{\nabla}$ and curvature $\overline{R}$, and
$M \subseteq \overline{M}$ is a submanifold with the induced metric $g = \bar{g}|_{TM}$, the \emph{first fundamental form}; we write $\langle -, - \rangle$ for
both. For $p \in M$, $T_p\overline{M} = T_pM \oplus N_pM$ with $N_pM = (T_pM)^\perp$, the \emph{normal space}, and $v = v^\top + v^\perp$ is the corresponding
decomposition. The normal spaces form the \emph{normal bundle} $NM$, a vector bundle over $M$: near each point it has a smooth frame, obtained by applying
Gram--Schmidt to a frame of $T\overline{M}$ extending one of $TM$. For $\overline{M} = \RR^N$ with the Euclidean metric, $\overline{\nabla}_XY$ is the directional
derivative of the components of $Y$, and $\overline{R} = 0$.

\begin{lemma}\label{riemannian-geometry:lemma-ambient-derivative}
Every vector field on an open subset of $M$ extends, near each point, to a vector field on an open subset of $\overline{M}$. For vector fields $X$, $Y$ on $M$ and
extensions $\overline{X}$, $\overline{Y}$, the value of $\overline{\nabla}_{\overline{X}}\overline{Y}$ at $p \in M$ depends only on $X$ and $Y$, and we denote it by
$\overline{\nabla}_XY$. Moreover $[\overline{X}, \overline{Y}]$ restricts on $M$ to $[X, Y]$.
\end{lemma}

\begin{proof}
In a slice chart $(x^1, \ldots, x^k, y^1, \ldots, y^m)$ of $\overline{M}$, in which $M$ is $\{y = 0\}$, a vector field $\sum_iY^i(x)\partial/\partial x^i$ on $M$ extends by
the same formula, with coefficients independent of $y$. The value of $\overline{\nabla}_{\overline{X}}\overline{Y}$ at $p$ depends only on $\overline{X}_p = X_p$, since a
connection is $C^\infty$-linear in its first argument, and it equals $D_t(\overline{Y} \circ \gamma)(0)$ for any curve $\gamma$ with $\dot\gamma(0) = X_p$
(Proposition~\ref{riemannian-geometry:proposition-covariant-derivative-curves}); taking $\gamma$ in $M$, $\overline{Y} \circ \gamma = Y \circ \gamma$. Finally $X$ and
$\overline{X}$ are related by the inclusion $\iota$, and so are $Y$ and $\overline{Y}$, so $[X, Y]$ and $[\overline{X}, \overline{Y}]$ are $\iota$-related
(Lemma~\ref{lie-groups:lemma-related-fields}).
\end{proof}

\begin{theorem}[Gauss formula]\label{riemannian-geometry:theorem-gauss-formula}
For vector fields $X$, $Y$ tangent to $M$,
\[
\overline{\nabla}_XY = \nabla_XY + \mathrm{II}(X, Y),
\]
where $\nabla_XY = (\overline{\nabla}_XY)^\top$ is the Levi-Civita connection of $(M, g)$, and $\mathrm{II}(X, Y) = (\overline{\nabla}_XY)^\perp$, the \emph{second fundamental
form}, is a symmetric $C^\infty(M)$-bilinear map with values in the normal bundle.
\end{theorem}

\begin{proof}
The operator $\nabla^\top_XY = (\overline{\nabla}_XY)^\top$ is a connection on $TM$: it is $C^\infty$-linear in $X$, and
$\overline{\nabla}_X(fY) = (Xf)Y + f\overline{\nabla}_XY$ with $Y$ tangent gives the Leibniz rule. It is metric: for tangent $Y$ and $Z$,
$X\langle Y, Z\rangle = \langle\overline{\nabla}_XY, Z\rangle + \langle Y, \overline{\nabla}_XZ\rangle = \langle\nabla^\top_XY, Z\rangle + \langle Y, \nabla^\top_XZ\rangle$, as the normal parts are orthogonal to
$Y$ and $Z$. By Lemma~\ref{riemannian-geometry:lemma-ambient-derivative} and the torsion-freeness of $\overline{\nabla}$,
$\overline{\nabla}_XY - \overline{\nabla}_YX = [X, Y]$, which is tangent. Its tangent part says that $\nabla^\top$ is torsion-free, so $\nabla^\top = \nabla$ by the uniqueness in
Theorem~\ref{riemannian-geometry:theorem-levi-civita}; its normal part says $\mathrm{II}(X, Y) = \mathrm{II}(Y, X)$. Finally $\mathrm{II}$ is $C^\infty$-linear in $X$, and in $Y$
because $\mathrm{II}(X, fY) = ((Xf)Y + f\overline{\nabla}_XY)^\perp = f\,\mathrm{II}(X, Y)$.
\end{proof}

So the intrinsic connection of $M$ is the tangential part of the ambient one: for a surface in $\RR^3$, $\nabla_XY$ is the derivative of $Y$ in $\RR^3$ projected
back onto the tangent plane. A curve in $M$ is a geodesic of $M$ exactly when its acceleration in $\overline{M}$ is normal to $M$; for instance great circles
on the sphere (Exercise~\ref{riemannian-geometry:exercise-great-circles}).

\begin{definition}\label{riemannian-geometry:definition-shape-operator}
For a normal vector field $\nu$ along $M$, the \emph{shape operator} (or \emph{Weingarten map}) is $S_\nu(X) = -(\overline{\nabla}_X\nu)^\top$, and the \emph{normal
connection} is $\nabla^\perp_X\nu = (\overline{\nabla}_X\nu)^\perp$.
\end{definition}

\begin{proposition}[Weingarten equation]\label{riemannian-geometry:proposition-weingarten}
For tangent $X$, $Y$ and normal $\nu$, $\langle S_\nu X, Y\rangle = \langle\mathrm{II}(X, Y), \nu\rangle$. In particular $S_\nu$ is self-adjoint on each tangent space, and
$\overline{\nabla}_X\nu = -S_\nu X + \nabla^\perp_X\nu$.
\end{proposition}

\begin{proof}
Differentiate $\langle Y, \nu\rangle = 0$: $0 = \langle\overline{\nabla}_XY, \nu\rangle + \langle Y, \overline{\nabla}_X\nu\rangle = \langle\mathrm{II}(X, Y), \nu\rangle - \langle Y, S_\nu X\rangle$. Self-adjointness
follows from the symmetry of $\mathrm{II}$, and the last formula is the decomposition of $\overline{\nabla}_X\nu$ into its two parts.
\end{proof}

\begin{theorem}[Gauss and Codazzi equations]\label{riemannian-geometry:theorem-gauss-codazzi}
For vector fields $X, Y, Z, W$ tangent to $M$:
\begin{enumerate}
\item (Gauss) $R(X, Y, Z, W) = \overline{R}(X, Y, Z, W) + \langle\mathrm{II}(Y, Z), \mathrm{II}(X, W)\rangle - \langle\mathrm{II}(X, Z), \mathrm{II}(Y, W)\rangle$;
\item (Codazzi) $(\overline{R}(X, Y)Z)^\perp = (\nabla_X\mathrm{II})(Y, Z) - (\nabla_Y\mathrm{II})(X, Z)$, where
$(\nabla_X\mathrm{II})(Y, Z) = \nabla^\perp_X(\mathrm{II}(Y, Z)) - \mathrm{II}(\nabla_XY, Z) - \mathrm{II}(Y, \nabla_XZ)$.
\end{enumerate}
In particular, for orthonormal $u, v \in T_pM$ the sectional curvatures of the plane they span satisfy
$K(u, v) = \overline{K}(u, v) + \langle\mathrm{II}(u, u), \mathrm{II}(v, v)\rangle - |\mathrm{II}(u, v)|^2$.
\end{theorem}

\begin{proof}
By the Gauss formula, $\overline{\nabla}_X\overline{\nabla}_YZ = \overline{\nabla}_X(\nabla_YZ) + \overline{\nabla}_X(\mathrm{II}(Y, Z))$, and applying the Gauss formula again to the tangent field
$\nabla_YZ$ and the Weingarten equation to the normal field $\mathrm{II}(Y, Z)$,
\[
\overline{\nabla}_X\overline{\nabla}_YZ = \nabla_X\nabla_YZ + \mathrm{II}(X, \nabla_YZ) - S_{\mathrm{II}(Y, Z)}X + \nabla^\perp_X(\mathrm{II}(Y, Z)).
\]
Also $\overline{\nabla}_{[X,Y]}Z = \nabla_{[X,Y]}Z + \mathrm{II}([X, Y], Z)$. Subtract the same expressions with $X$ and $Y$ exchanged, and use $[X, Y] = \nabla_XY - \nabla_YX$.
The tangential part gives $(\overline{R}(X, Y)Z)^\top = R(X, Y)Z - S_{\mathrm{II}(Y, Z)}X + S_{\mathrm{II}(X, Z)}Y$; pairing with $W$ and using the Weingarten equation gives (1). The
normal part gives
\[
(\overline{R}(X, Y)Z)^\perp = \nabla^\perp_X(\mathrm{II}(Y, Z)) + \mathrm{II}(X, \nabla_YZ) - \nabla^\perp_Y(\mathrm{II}(X, Z)) - \mathrm{II}(Y, \nabla_XZ) - \mathrm{II}(\nabla_XY - \nabla_YX, Z),
\]
which is (2) after regrouping. For the last statement take $X = W = u$ and $Y = Z = v$ in (1), recalling $K(u, v) = R(u, v, v, u)$.
\end{proof}

\begin{definition}\label{riemannian-geometry:definition-hypersurface-curvatures}
Let $M$ be a hypersurface ($\dim\overline{M} = \dim M + 1$) with a unit normal field $\nu$, which exists locally, and globally if $M$ and $\overline{M}$ are orientable.
The \emph{scalar second fundamental form} is $h(X, Y) = \langle\mathrm{II}(X, Y), \nu\rangle$, so $\mathrm{II} = h\,\nu$, and the shape operator $S = S_\nu$ satisfies
$\langle SX, Y\rangle = h(X, Y)$. The eigenvalues $\kappa_1, \ldots, \kappa_n$ of the self-adjoint map $S$ on $T_pM$ are the \emph{principal curvatures}, its eigenvectors the
\emph{principal directions}, $H = \frac1n\operatorname{tr}S$ is the \emph{mean curvature}, and $\det S = \kappa_1 \cdots \kappa_n$ the \emph{Gauss--Kronecker
curvature}. Changing $\nu$ to $-\nu$ changes the signs of $h$, $S$, the $\kappa_i$ and $H$.
\end{definition}

\begin{corollary}[Theorema Egregium]\label{riemannian-geometry:corollary-theorema-egregium}
Let $M^n \subseteq \RR^{n+1}$ be a hypersurface with unit normal $\nu$. Then
$R(X, Y, Z, W) = h(Y, Z)h(X, W) - h(X, Z)h(Y, W)$, and for orthonormal principal directions $e_i$, $e_j$ the sectional curvature is $K(e_i, e_j) = \kappa_i\kappa_j$. In
particular, for a surface in $\RR^3$ the Gaussian curvature is $K = \kappa_1\kappa_2 = \det S$: a quantity defined through the embedding is determined by the metric
alone, and so is preserved by isometries between surfaces.
\end{corollary}

\begin{proof}
Apply Theorem~\ref{riemannian-geometry:theorem-gauss-codazzi}(1) with $\overline{R} = 0$ and $\mathrm{II} = h\nu$, $|\nu| = 1$. For principal directions,
$h(e_i, e_i) = \kappa_i$ and $h(e_i, e_j) = 0$ for $i \neq j$. For a surface the only plane is the tangent plane, and $\kappa_1\kappa_2$ is the determinant of $S$.
\end{proof}

Codazzi's equation for a hypersurface of $\RR^{n+1}$ says that $(\nabla_Xh)(Y, Z)$ is symmetric in all three arguments. Gauss's and Codazzi's equations are the
integrability conditions for the embedding: by Bonnet's theorem, a metric and a symmetric form on a simply connected surface satisfying them come from a surface
in $\RR^3$, unique up to rigid motions \cite[Chapter 4]{do-carmo-curves-surfaces}.

\begin{proposition}[Surfaces in coordinates]\label{riemannian-geometry:proposition-surface-coordinates}
Let $x \colon \Omega \to \RR^3$ be a parametrisation of a surface $S$ (Remark~\ref{smooth-manifolds:remark-do-carmo}), with coordinates $(u^1, u^2) = (u, v)$,
$x_i = \partial x/\partial u^i$, and unit normal $\nu = x_1 \times x_2/|x_1 \times x_2|$. Let
\[
g_{ij} = \langle x_i, x_j\rangle = \begin{pmatrix} E & F\\ F & G\end{pmatrix}, \qquad h_{ij} = \langle\partial_i\partial_jx, \nu\rangle = -\langle x_i, \partial_j\nu\rangle.
\]
Then the first and second fundamental forms are $E\,du^2 + 2F\,du\,dv + G\,dv^2$ and $h_{11}\,du^2 + 2h_{12}\,du\,dv + h_{22}\,dv^2$, the matrix of $S$ in the basis
$(x_1, x_2)$ is $(g_{ij})^{-1}(h_{ij})$, and
\[
K = \frac{h_{11}h_{22} - h_{12}^2}{EG - F^2}, \qquad H = \frac{Eh_{22} - 2Fh_{12} + Gh_{11}}{2(EG - F^2)}.
\]
More generally, for a submanifold of $\RR^N$ parametrised by $x(u^1, \ldots, u^n)$, the decomposition
$\partial_i\partial_jx = \Gamma^k_{ij}x_k + \mathrm{II}(x_i, x_j)$ of the second derivatives into tangent and normal parts gives the Christoffel symbols and the second
fundamental form, and $R_{ijkl} = R(x_i, x_j, x_k, x_l) = \langle\mathrm{II}_{jk}, \mathrm{II}_{il}\rangle - \langle\mathrm{II}_{ik}, \mathrm{II}_{jl}\rangle$ with
$\mathrm{II}_{ij} = \mathrm{II}(x_i, x_j)$.
\end{proposition}

\begin{proof}
In $\RR^N$, $\overline{\nabla}$ is the directional derivative, and the coordinate field $\partial/\partial u^j$ is $x_j$ along $S$, so
$\overline{\nabla}_{x_i}x_j = \partial_i\partial_jx$. The Gauss formula splits it as $\nabla_{x_i}x_j + \mathrm{II}(x_i, x_j)$, which gives the general statement, the tangent part
being $\Gamma^k_{ij}x_k$, and for surfaces $h_{ij} = \langle\partial_i\partial_jx, \nu\rangle$; differentiating $\langle x_i, \nu\rangle = 0$ gives the second expression. From
$\langle Sx_i, x_j\rangle = h_{ij}$, the matrix $A$ of $S$ satisfies $(g_{ij})A = (h_{ij})$. So $K = \det S = \det(h_{ij})/\det(g_{ij})$ and
$2H = \operatorname{tr}((g_{ij})^{-1}(h_{ij}))$, which are the formulas stated, using $(g_{ij})^{-1} = \frac{1}{EG - F^2}\bigl(\begin{smallmatrix} G & -F\\ -F & E\end{smallmatrix}\bigr)$. The
formula for $R_{ijkl}$ is Theorem~\ref{riemannian-geometry:theorem-gauss-codazzi}(1) with $\overline{R} = 0$.
\end{proof}

\begin{example}\label{riemannian-geometry:example-surfaces}
\begin{enumerate}
\item \emph{Spheres.} For the sphere of radius $r$ in $\RR^{n+1}$ with outward normal $\nu(p) = p/r$, $\overline{\nabla}_X\nu = X/r$, so $S = -\frac1r\id$, $h = -\frac1rg$, every
direction is principal with $\kappa = -1/r$, and all sectional curvatures are $\kappa_i\kappa_j = 1/r^2$, as in Example~\ref{riemannian-geometry:example-curvature}.
\item \emph{Cylinders.} For $\{x^2 + y^2 = 1\} \subseteq \RR^3$ with $\nu = (x, y, 0)$, $\overline{\nabla}_X\nu$ is the horizontal part of $X$, so the principal curvatures are
$-1$ (horizontal circles) and $0$ (vertical lines), $K = 0$ and $H = -\frac12$. The cylinder is locally isometric to the plane (unroll it), in accordance with the Theorema
Egregium, although it bends in $\RR^3$.
\item \emph{Graphs.} Let $S$ be the graph of $z = \phi(u, v)$ with $\phi(0) = 0$ and $d\phi(0) = 0$, parametrised by $x(u, v) = (u, v, \phi(u, v))$. At the origin
$x_1 = e_1$, $x_2 = e_2$, $\nu = e_3$, so $g_{ij} = \delta_{ij}$ and $h_{ij} = \partial_i\partial_j\phi(0)$: the second fundamental form is the Hessian of $\phi$, the principal
curvatures are its eigenvalues, and $K(0) = \det\operatorname{Hess}\phi(0)$. The paraboloid $z = u^2 + v^2$ has $K(0) = 4$, the saddle $z = u^2 - v^2$ has $K(0) = -4$: positive
curvature means that the surface lies locally on one side of its tangent plane, negative curvature that it crosses it.
\end{enumerate}
\end{example}

\section{Global theorems on surfaces}\label{riemannian-geometry:section-surfaces-global}

This section collects the classical global theorems on surfaces, in the form in which they appear in \cite[Chapters 4 and 5]{do-carmo-curves-surfaces}:
the Gauss--Bonnet theorem for regions with boundary, Minding's theorem on surfaces of constant curvature, and the theorems of Hadamard, Liebmann and Hilbert on
surfaces in $\RR^3$. Throughout, $M$ is an oriented Riemannian surface and $J$ is the rotation by $+\pi/2$ in its tangent planes, as in the proof of
Theorem~\ref{riemannian-geometry:theorem-gauss-bonnet}; for a surface in $\RR^3$ oriented by a unit normal $\nu$, $Jv = \nu \times v$
(Proposition~\ref{differential-forms:proposition-orientation-normal}).

\begin{definition}\label{riemannian-geometry:definition-geodesic-curvature}
Let $\alpha$ be a curve in $M$ parametrised by arc length, $|\alpha'| = 1$. Since $\langle D_s\alpha', \alpha'\rangle = \frac12\frac{d}{ds}|\alpha'|^2 = 0$, the covariant acceleration is
$D_s\alpha' = k_g\,J\alpha'$ for a function $k_g$, the \emph{geodesic curvature} of $\alpha$. So $\alpha$ is a geodesic if and only if $k_g = 0$. For a surface in $\RR^3$,
$D_s\alpha'$ is the tangential part of the acceleration $\alpha''$ (Theorem~\ref{riemannian-geometry:theorem-gauss-formula}), and $k_g = \langle\alpha'', \nu \times \alpha'\rangle$.
\end{definition}

For example, a circle of latitude at height $z_0 \in (-1, 1)$ on the unit sphere has geodesic curvature $\pm z_0/(1 - z_0^2)^{1/2}$, and the equator, $z_0 = 0$, is a geodesic.

A \emph{simple region} is a compact set $R \subseteq M$ contained in the domain of a positively oriented chart $x \colon \Omega \to M$, whose preimage is homeomorphic to a closed
disc and bounded by a simple closed curve that is piecewise smooth with nonvanishing velocity. Its boundary $\partial R$ is parametrised by arc length so that $R$ lies on the left
(that is, $J\alpha'$ points into $R$), and at a vertex where the smooth arcs meet, the \emph{exterior angle} $\theta_i \in (-\pi, \pi)$ is the oriented angle from the incoming to the
outgoing velocity; the \emph{interior angle} is $\pi - \theta_i$.

\begin{theorem}[Theorem of turning tangents]\label{riemannian-geometry:theorem-turning-tangents}
Let $R$ be a simple region, $(e_1, e_2)$ a positively oriented orthonormal frame on a neighbourhood of $R$, and write the unit tangent of $\partial R$ on each smooth arc as
$\alpha' = \cos\varphi\,e_1 + \sin\varphi\,e_2$ with $\varphi$ continuous on the arc. Then the sum over the arcs of the changes of $\varphi$, plus the sum of the exterior angles,
equals $2\pi$.
\end{theorem}

This is Hopf's Umlaufsatz \cite{hopf-1935}; see \cite[Chapters 4 and 5]{do-carmo-curves-surfaces}. For a smooth convex curve in the Euclidean plane it says that the tangent turns
exactly once.

\begin{theorem}[Local Gauss--Bonnet theorem]\label{riemannian-geometry:theorem-local-gauss-bonnet}
Let $R$ be a simple region with smooth boundary arcs $\alpha_i$ and exterior angles $\theta_1, \ldots, \theta_k$. Then
\[
\sum_i\int_{\alpha_i}k_g\,ds + \iint_RK\,dA + \sum_{j=1}^k\theta_j = 2\pi.
\]
\end{theorem}

\begin{proof}
Applying Gram--Schmidt to the coordinate fields of the chart gives a positively oriented orthonormal frame $(e_1, e_2)$ near $R$. Define the $1$-form $\omega$ by
$\nabla_Xe_1 = \omega(X)e_2$; then $\nabla_Xe_2 = -\omega(X)e_1$, and the computation in the proof of Theorem~\ref{riemannian-geometry:theorem-gauss-bonnet} gives
$d\omega = -K\,dA$. On a smooth arc, $\alpha' = \cos\varphi\,e_1 + \sin\varphi\,e_2$ with $\varphi$ smooth (lift the map to the circle through the covering $\RR \to S^1$,
Proposition~\ref{homotopy:proposition-homotopy-lifting}), and $J\alpha' = -\sin\varphi\,e_1 + \cos\varphi\,e_2$. Differentiating,
\[
D_s\alpha' = \varphi'J\alpha' + \cos\varphi\,\nabla_{\alpha'}e_1 + \sin\varphi\,\nabla_{\alpha'}e_2 = \varphi'J\alpha' + \omega(\alpha')J\alpha',
\]
so $k_g = \varphi' + \omega(\alpha')$. Summing over the arcs, $\sum_i\int_{\alpha_i}k_g\,ds = \sum_i(\text{change of }\varphi\text{ on }\alpha_i) + \int_{\partial R}\omega$. The first sum is
$2\pi - \sum_j\theta_j$ by Theorem~\ref{riemannian-geometry:theorem-turning-tangents}, and $\int_{\partial R}\omega = \iint_Rd\omega = -\iint_RK\,dA$ by Stokes' theorem in the chart, that
is, Green's theorem for a plane region bounded by a piecewise smooth simple closed curve \cite[Chapter 4]{do-carmo-curves-surfaces}
(Theorem~\ref{differential-forms:theorem-stokes} in the smooth case).
\end{proof}

\begin{corollary}\label{riemannian-geometry:corollary-geodesic-triangles}
Let $T$ be a \emph{geodesic triangle}: a simple region bounded by three geodesic arcs, with interior angles $\phi_1, \phi_2, \phi_3$. Then
\[
\phi_1 + \phi_2 + \phi_3 - \pi = \iint_TK\,dA.
\]
So the angles of a geodesic triangle add up to more than $\pi$ where $K > 0$, to $\pi$ where $K = 0$, and to less than $\pi$ where $K < 0$. On the unit sphere the area of a geodesic
triangle is its angle excess; for instance the octant bounded by the three coordinate planes has three right angles and area $\frac{3\pi}{2} - \pi = \frac{\pi}{2}$, one eighth of $4\pi$.
\end{corollary}

\begin{proof}
Here $k_g = 0$ on the sides and $\theta_j = \pi - \phi_j$, so Theorem~\ref{riemannian-geometry:theorem-local-gauss-bonnet} gives
$\iint_TK\,dA + 3\pi - (\phi_1 + \phi_2 + \phi_3) = 2\pi$.
\end{proof}

\begin{theorem}[Global Gauss--Bonnet theorem]\label{riemannian-geometry:theorem-global-gauss-bonnet}
Let $R \subseteq M$ be a compact region whose boundary consists of finitely many simple closed piecewise smooth curves $C_1, \ldots, C_m$, oriented with $R$ on the left, with
exterior angles $\theta_1, \ldots, \theta_k$ at their vertices. Then
\[
\iint_RK\,dA + \sum_i\int_{C_i}k_g\,ds + \sum_j\theta_j = 2\pi\chi(R).
\]
\end{theorem}

\begin{proof}[Sketch of proof]
Choose a triangulation of $R$ fine enough that every triangle is a simple region, with $F$ triangles, $E$ edges and $V$ vertices, so that $\chi(R) = V - E + F$. Apply
Theorem~\ref{riemannian-geometry:theorem-local-gauss-bonnet} to each triangle and add. Every interior edge is traversed twice in opposite directions, so the integrals of $k_g$ over
interior edges cancel, and those over boundary edges give the boundary term. The exterior angles of a triangle are $\pi$ minus its interior angles, and the interior
angles of all triangles add up to $2\pi$ at each of the $V_{\mathrm{int}}$ interior vertices, to $\pi - \theta_j$ at a corner of $\partial R$, and to $\pi$ at the other boundary
vertices. So the sum of all exterior angles of the triangles is $3\pi F - 2\pi V_{\mathrm{int}} - \pi V_\partial + \sum_j\theta_j$, where $V_\partial$ is the number of boundary
vertices, and adding the local formulas gives
$\iint_RK\,dA + \sum_i\int_{C_i}k_g\,ds + \sum_j\theta_j = -\pi F + 2\pi V_{\mathrm{int}} + \pi V_\partial$. Each triangle has three edges, each interior edge lies in two triangles
and each boundary edge in one, so $3F = 2E_{\mathrm{int}} + E_\partial$, and $E_\partial = V_\partial$ because the boundary is a union of closed curves; with these relations the right
side equals $2\pi(V - E + F)$. The omitted verification is the existence of such triangulations \cite[Chapter 4]{do-carmo-curves-surfaces}.
\end{proof}

For $R = M$ a closed surface there is no boundary, and the formula is Theorem~\ref{riemannian-geometry:theorem-gauss-bonnet}, $\iint_MK\,dA = 2\pi\chi(M)$. For a disc in the
plane bounded by a smooth convex curve, $K = 0$ and $\chi = 1$, and the formula says that the total geodesic curvature of the boundary is $2\pi$: the tangent turns once. For the
northern cap of the unit sphere above the latitude $z_0$, the area $2\pi(1 - z_0)$ plus the total geodesic curvature $2\pi(1 - z_0^2)^{1/2} \cdot z_0/(1 - z_0^2)^{1/2} = 2\pi z_0$ of
the boundary circle gives $2\pi = 2\pi\chi(\text{disc})$.

\begin{theorem}[Minding; Cartan]\label{riemannian-geometry:theorem-minding}
Let $(M, g)$ and $(N, h)$ be Riemannian manifolds of the same dimension $n$ with the same constant sectional curvature $\kappa$, and let $p \in M$, $q \in N$ and a linear isometry
$I \colon T_pM \to T_qN$ be given. Then there are neighbourhoods $U \ni p$ and $V \ni q$ and an isometry $F \colon U \to V$ with $F(p) = q$ and $d_pF = I$. In particular two surfaces
with the same constant Gaussian curvature are locally isometric, and every surface of constant curvature $0$, $1$ or $-1$ is locally isometric to the plane, the unit sphere or the
hyperbolic plane.
\end{theorem}

\begin{proof}
Let $s_\kappa$ be the solution of $s'' + \kappa s = 0$ with $s(0) = 0$ and $s'(0) = 1$ (so $s_\kappa(t) = t$, $\sin(\sqrt\kappa t)/\sqrt\kappa$ or $\sinh(\sqrt{-\kappa}t)/\sqrt{-\kappa}$). We
first compute the differential of $\exp_p$. Let $v \in T_pM$ be in the domain of $\exp_p$, $v \neq 0$, $\ell = |v|$, $u = v/\ell$, and $\gamma(t) = \exp_p(tv)$; for $w \in T_pM$ write
$w = w_\parallel + w_\perp$ with $w_\parallel$ a multiple of $u$ and $w_\perp \perp u$, and let $P_t$ be parallel transport along $\gamma$. By
Proposition~\ref{riemannian-geometry:proposition-jacobi}(2), $J(t) = d_{tv}\exp_p(tw)$ is the Jacobi field along $\gamma$ with $J(0) = 0$ and $D_tJ(0) = w$. The fields
$J_1(t) = tP_t(w_\parallel)$ and $J_2(t) = \ell^{-1}s_\kappa(\ell t)P_t(w_\perp)$ are Jacobi fields: $P_t(u) = \dot\gamma(t)/\ell$, so $J_1$ is a multiple of $t\dot\gamma$, with $D_t^2J_1 = 0$ and
$R(J_1, \dot\gamma)\dot\gamma = 0$; and $P_t(w_\perp) \perp \dot\gamma$, so by the curvature formula of Example~\ref{riemannian-geometry:example-curvature},
$R(J_2, \dot\gamma)\dot\gamma = \kappa|\dot\gamma|^2J_2 = \kappa\ell^2J_2$, while $D_t^2J_2 = \ell s_\kappa''(\ell t)P_t(w_\perp) = -\kappa\ell^2J_2$. They have the initial values $J_1(0) = J_2(0) = 0$,
$D_tJ_1(0) = w_\parallel$ and $D_tJ_2(0) = w_\perp$, so $J = J_1 + J_2$ by uniqueness, and since $P_1$ is an isometry and $P_1(w_\parallel) \perp P_1(w_\perp)$,
\[
|d_v\exp_p(w)|^2 = |J(1)|^2 = |w_\parallel|^2 + \Bigl(\frac{s_\kappa(\ell)}{\ell}\Bigr)^2|w_\perp|^2.
\]
The same formula holds on $N$. Now choose $\varepsilon > 0$ such that $\exp_p$ and $\exp_q$ map the balls of radius $\varepsilon$ diffeomorphically onto open sets $U$ and $V$
(Proposition~\ref{riemannian-geometry:proposition-exponential}), and let $F = \exp_q \circ I \circ (\exp_p|_{B_\varepsilon})^{-1} \colon U \to V$, a diffeomorphism with $F(p) = q$ and $d_pF = I$.
For $x = \exp_p(v)$ and $w \in T_pM$, the chain rule gives $d_xF(d_v\exp_p(w)) = d_{Iv}\exp_q(Iw)$, and since $I$ preserves lengths and orthogonality, the formula above gives the same
length for $d_v\exp_p(w)$ and $d_{Iv}\exp_q(Iw)$ (for $v = 0$ both differentials are the identity). As the $d_v\exp_p(w)$ fill $T_xM$, $d_xF$ preserves lengths, hence inner products, so $F$
is an isometry.
\end{proof}

For instance the cylinder and the cone (without its vertex) are locally isometric to the plane, and the pseudosphere, the surface of revolution of the tractrix, to the hyperbolic
plane. Minding's theorem is local: the flat torus $\RR^2/\ZZ^2$ is locally but not globally isometric to the plane.

\begin{lemma}\label{riemannian-geometry:lemma-elliptic-point}
A compact surface $S \subseteq \RR^3$ has a point where $K > 0$. More precisely, if $p_0 \in S$ maximises the distance $R$ to the origin, then $K(p_0) \geq 1/R^2$.
\end{lemma}

\begin{proof}
The function $|x|^2$ on $S$ has a maximum $R^2 > 0$ at some $p_0$, since $S$ is compact (and not a point). Its differential $2\langle p_0, -\rangle$ vanishes on $T_{p_0}S$, so
$\nu = p_0/R$ is a unit normal at $p_0$. For a unit-speed curve $\alpha$ in $S$ with $\alpha(0) = p_0$, the function $|\alpha|^2$ has a maximum at $0$, so
$0 \geq (|\alpha|^2)''(0) = 2\langle\alpha''(0), p_0\rangle + 2|\alpha'(0)|^2$, that is, $\langle\alpha''(0), \nu\rangle \leq -1/R$. The normal part of $\alpha''(0)$ is
$\mathrm{II}(\alpha'(0), \alpha'(0))$ (Theorem~\ref{riemannian-geometry:theorem-gauss-formula}, applied to the covariant acceleration of $\alpha$ in $\RR^3$), so
$h(w, w) \leq -1/R$ for every unit tangent vector $w$ at $p_0$. Both principal curvatures are therefore $\leq -1/R$, and $K(p_0) = \kappa_1\kappa_2 \geq 1/R^2$
(Corollary~\ref{riemannian-geometry:corollary-theorema-egregium}).
\end{proof}

So there is no compact surface in $\RR^3$ with $K \leq 0$ everywhere; in particular the flat torus cannot be embedded isometrically in $\RR^3$ as a smooth surface.

\begin{lemma}\label{riemannian-geometry:lemma-proper-covering}
Let $f \colon X \to Y$ be a local homeomorphism from a compact Hausdorff space to a connected Hausdorff space. Then $f$ is surjective and a covering map with finite fibres.
\end{lemma}

\begin{proof}
A local homeomorphism is open, and a continuous map from a compact space to a Hausdorff space is closed (Proposition~\ref{general-topology:proposition-compact-basic}); so $f(X)$ is
open, closed and nonempty, hence equal to $Y$. A fibre $f^{-1}(y)$ is discrete, because $f$ is injective near each point, and closed in the compact $X$, hence finite, say
$\{x_1, \ldots, x_k\}$. Choose pairwise disjoint open sets $U_i \ni x_i$ mapped homeomorphically onto open sets $V_i \ni y$. The set $f(X \setminus \bigcup_iU_i)$ is compact, hence
closed, and does not contain $y$, so $V = \bigcap_iV_i \setminus f(X \setminus \bigcup_iU_i)$ is an open neighbourhood of $y$ with $f^{-1}(V) = \bigsqcup_i(U_i \cap f^{-1}(V))$, each piece
mapped homeomorphically onto $V$.
\end{proof}

\begin{theorem}[Hadamard]\label{riemannian-geometry:theorem-hadamard}
Let $S \subseteq \RR^3$ be a compact connected surface with Gaussian curvature $K > 0$ everywhere. Then $S$ is orientable, its Gauss map $\nu \colon S \to S^2$ (for the unit normal
making the second fundamental form negative definite) is a diffeomorphism, $S$ is diffeomorphic to $S^2$, and $\iint_SK\,dA = 4\pi$.
\end{theorem}

\begin{proof}
At each point $K = \det S_\nu > 0$, so the self-adjoint shape operator is definite, and exactly one of the two unit normals makes it negative definite; near each point this normal is
$\pm$ a smooth local normal with a locally constant sign, so it is a unit normal field, and $S$ is orientable (Proposition~\ref{differential-forms:proposition-orientation-normal}). For
tangent $X$, $\overline{\nabla}_X\nu$ is tangent, because $\langle\nu, \nu\rangle = 1$, so $d_p\nu(X) = \overline{\nabla}_X\nu = -S_\nu X$
(Proposition~\ref{riemannian-geometry:proposition-weingarten}); as $T_{\nu(p)}S^2 = \nu(p)^\perp = T_pS$ and $\det(-S_\nu) = K > 0$, $d_p\nu$ is an isomorphism, and $\nu$ is a local
diffeomorphism (Theorem~\ref{smooth-manifolds:theorem-inverse-function-manifolds}). By Lemma~\ref{riemannian-geometry:lemma-proper-covering} it is a covering map, and since $S^2$ is simply
connected (Corollary~\ref{homotopy:corollary-spheres}) and $S$ connected, it has one sheet (Corollary~\ref{homotopy:corollary-lifting-consequences}(3)): $\nu$ is a bijective local
diffeomorphism, hence a diffeomorphism. Orient $S$ by $\nu$ and $S^2$ by its outward normal; then $T_pS$ and $T_{\nu(p)}S^2$ are the same oriented plane, $\nu^*(dA_{S^2}) = \det(d\nu)\,dA_S = K\,dA_S$,
and $\iint_SK\,dA = \int_{S^2}dA = 4\pi$ (Proposition~\ref{differential-forms:proposition-integral-well-defined}).
\end{proof}

The value $4\pi$ agrees with the Gauss--Bonnet theorem, $2\pi\chi(S^2)$. Hadamard also proved that such a surface is the boundary of a convex body: it lies on one side of each
of its tangent planes \cite[Chapter 5]{do-carmo-curves-surfaces}.

\begin{theorem}[Liebmann]\label{riemannian-geometry:theorem-liebmann}
A compact connected surface in $\RR^3$ with constant Gaussian curvature $K$ is a round sphere of radius $1/\sqrt{K}$.
\end{theorem}

\begin{proof}[Sketch of proof]
By Lemma~\ref{riemannian-geometry:lemma-elliptic-point} the constant $K$ is positive. With the unit normal $\nu$ of Theorem~\ref{riemannian-geometry:theorem-hadamard} the
principal curvatures are negative; write them as $-\lambda_1 \leq -\lambda_2$, so that $\lambda_1 \geq \lambda_2 > 0$ and $\lambda_1\lambda_2 = K$. The continuous function $\lambda_1$
attains its maximum at some $p$, where $\lambda_2 = K/\lambda_1$ attains its minimum. \emph{Hilbert's lemma}, a computation with the Codazzi equation in principal coordinates, says
that at a non-umbilic point ($\lambda_1 > \lambda_2$) where $\lambda_1$ has a local maximum and $\lambda_2$ a local minimum, $K \leq 0$. So $p$ is umbilic and
$\lambda_1(p) = \sqrt K$. Hence $\lambda_1 \leq \sqrt K$ everywhere, while $\lambda_1^2 \geq \lambda_1\lambda_2 = K$; so $\lambda_1 = \lambda_2 = \sqrt K$ everywhere, every point is umbilic,
and $S_\nu = -\sqrt K\,\id$. Then $d\nu = -S_\nu = \sqrt K\,\id$, so $x - \nu/\sqrt K$ is constant on the connected surface $S$: every point lies at distance $1/\sqrt K$ from this
fixed point. The omitted verification is Hilbert's lemma \cite[Chapter 5]{do-carmo-curves-surfaces}.
\end{proof}

\begin{theorem}[Hilbert]\label{riemannian-geometry:theorem-hilbert}
There is no complete surface of constant negative Gaussian curvature isometrically immersed in $\RR^3$. In particular the hyperbolic plane cannot be realised as a complete
surface in $\RR^3$.
\end{theorem}

This is \cite{hilbert-1901}; see \cite[Chapter 5]{do-carmo-curves-surfaces} for a complete proof. In outline: on such a surface the asymptotic directions (where $h(w, w) = 0$) give
coordinates in which the metric has the form $du^2 + 2\cos\theta\,du\,dv + dv^2$, and the Gauss equation becomes the sine-Gordon equation $\theta_{uv} = \sin\theta$; these
coordinates cover the whole hyperbolic plane, whose area is infinite, while the Gauss--Bonnet theorem bounds the area of every coordinate quadrilateral by $2\pi$, a contradiction.
Incomplete surfaces of constant negative curvature do exist, such as the pseudosphere; and by Theorem~\ref{riemannian-geometry:theorem-minding} each is locally isometric to
the hyperbolic plane.

\section{The Hodge star and the Laplacian}\label{riemannian-geometry:section-hodge-star}

\begin{definition}\label{riemannian-geometry:definition-hodge-star}
Let $(M, g)$ be an oriented Riemannian $n$-manifold. The metric induces inner products on $\Lambda^kT^*_pM$ (Proposition~\ref{multilinear-algebra:proposition-induced-inner-product}), a volume form
$\mathrm{vol}_g$, and the \emph{Hodge star} $* \colon \Omega^k(M) \to \Omega^{n-k}(M)$, characterised pointwise by $\alpha \wedge *\beta = \langle\alpha, \beta\rangle\,\mathrm{vol}_g$
(Theorem~\ref{multilinear-algebra:theorem-hodge-star}); it satisfies $** = (-1)^{k(n-k)}$ on $k$-forms. In coordinates $\mathrm{vol}_g = \sqrt{\det(g_{ij})}\,dx^1 \wedge \cdots \wedge dx^n$ in a positively
oriented chart.
\end{definition}

\begin{definition}\label{riemannian-geometry:definition-laplacian}
On a compact oriented Riemannian manifold define the $L^2$ inner product on $\Omega^k(M)$ by
\[
(\alpha, \beta) = \int_M\alpha \wedge *\beta = \int_M\langle\alpha, \beta\rangle\,\mathrm{vol}_g,
\]
the \emph{codifferential} $\delta = (-1)^{n(k+1)+1}*d* \colon \Omega^k(M) \to \Omega^{k-1}(M)$, and the \emph{Hodge Laplacian} $\Delta = d\delta + \delta d$. A form is \emph{harmonic} if $\Delta\alpha = 0$; the space of
harmonic $k$-forms is $\mathcal{H}^k(M)$.
\end{definition}

\begin{proposition}\label{riemannian-geometry:proposition-adjoint}
On a compact oriented Riemannian manifold:
\begin{enumerate}
\item $(d\alpha, \beta) = (\alpha, \delta\beta)$ for $\alpha \in \Omega^{k-1}$, $\beta \in \Omega^k$: the codifferential is the formal adjoint of $d$;
\item $\Delta$ is self-adjoint and non-negative: $(\Delta\alpha, \alpha) = \|d\alpha\|^2 + \|\delta\alpha\|^2$;
\item $\Delta\alpha = 0$ if and only if $d\alpha = 0$ and $\delta\alpha = 0$;
\item $*\Delta = \Delta*$, so $*$ maps harmonic $k$-forms isomorphically to harmonic $(n-k)$-forms;
\item $\Delta$ on functions is $-\sum_{i}\partial_i^2$ in normal coordinates at a point, and in general
$\Delta f = -\frac{1}{\sqrt{\det g}}\sum_{i,j}\partial_i(\sqrt{\det g}\,g^{ij}\partial_jf)$.
\end{enumerate}
\end{proposition}

\begin{proof}
(1) By Stokes' theorem (Theorem~\ref{differential-forms:theorem-stokes}), $0 = \int_Md(\alpha \wedge *\beta) = \int_Md\alpha \wedge *\beta + (-1)^{k-1}\int_M\alpha \wedge d{*}\beta$. Using $** = (-1)^{k(n-k)}$ one
rewrites $d*\beta = \pm *{*}d{*}\beta$, and the sign in the definition of $\delta$ is chosen exactly so that the identity becomes $(d\alpha, \beta) = (\alpha, \delta\beta)$.

(2) $(\Delta\alpha, \beta) = (d\delta\alpha, \beta) + (\delta d\alpha, \beta) = (\delta\alpha, \delta\beta) + (d\alpha, d\beta)$, which is symmetric; taking $\beta = \alpha$ gives the formula. (3) If $\Delta\alpha = 0$ then
$\|d\alpha\|^2 + \|\delta\alpha\|^2 = 0$, so both vanish; the converse is clear. (4) follows from the definition of $\delta$ and $** = \pm 1$: $*d\delta = \delta d*$ and $*\delta d = d\delta*$ up to the signs, which
cancel. (5) is a computation with the coordinate formula for $*$ and $d$, using $\mathrm{vol}_g = \sqrt{\det g}\,dx^1\wedge\cdots\wedge dx^n$; in normal coordinates at $p$ the first derivatives of $g$ vanish, so
$\Delta f(p) = -\sum_i\partial_i^2f(p)$.
\end{proof}

\begin{remark}\label{riemannian-geometry:remark-weitzenbock}
The Levi-Civita connection gives a second Laplacian, the \emph{connection Laplacian} $\nabla^*\nabla$, and the two differ by curvature: the \emph{Weitzenb\"ock formula} states
$\Delta = \nabla^*\nabla + \mathcal{R}$, where $\mathcal{R}$ is an endomorphism built from the curvature tensor; on $1$-forms $\mathcal{R} = \mathrm{Ric}$. Since $(\nabla^*\nabla\alpha, \alpha) = \|\nabla\alpha\|^2 \geq 0$,
this gives vanishing theorems: if $\mathrm{Ric} > 0$ on a compact oriented manifold, every harmonic $1$-form vanishes, hence $b_1(M) = 0$ by the Hodge theorem (Bochner's theorem,
Chapter~\ref{hodge-theorem}). For proofs see \cite[Chapter 6]{warner-foundations} and \cite[Chapter 4]{lee-riemannian}.
\end{remark}

\begin{example}\label{riemannian-geometry:example-harmonic-forms}
On a compact connected oriented manifold the harmonic $0$-forms are the constants, since $\Delta f = 0$ implies $df = 0$; by (4) the harmonic $n$-forms are the multiples of $\mathrm{vol}_g$. On the flat torus
$\RR^n/\ZZ^n$, the harmonic forms are exactly the forms with constant coefficients in the coordinates $dx^i$, a space of dimension $\binom{n}{k}$, which matches $b_k(T^n)$
(Exercise~\ref{de-rham:exercise-torus}).
\end{example}

\section{Exercises}\label{riemannian-geometry:section-exercises}

\begin{exercise}\label{riemannian-geometry:exercise-hodge-star-signs}
Verify that $\delta = -*d*$ on $k$-forms when $n$ is even, and compute $\delta$ on $1$-forms in $\RR^3$, identifying it with minus the divergence.
\end{exercise}

\begin{exercise}\label{riemannian-geometry:exercise-conformal}
Show that in dimension $n = 2k$ the Hodge star on $k$-forms depends only on the conformal class of the metric, and that $*^2 = (-1)^k$ on $\Omega^k$.
\end{exercise}

\begin{solution}
Under $g \mapsto \lambda g$ with $\lambda > 0$, the inner product on $k$-forms scales by $\lambda^{-k}$ and the volume form by $\lambda^{n/2} = \lambda^k$, so $*$ is unchanged by the characterisation
$\alpha \wedge *\beta = \langle\alpha,\beta\rangle\mathrm{vol}$. The sign is $(-1)^{k(n-k)} = (-1)^{k^2} = (-1)^k$.
\end{solution}

\begin{exercise}\label{riemannian-geometry:exercise-laplacian-sphere}
Compute the Laplacian on functions on the round sphere $S^2$ in spherical coordinates, and verify that the restrictions of linear functions of $\RR^3$ are eigenfunctions with eigenvalue $2$.
\end{exercise}

\begin{exercise}\label{riemannian-geometry:exercise-bi-invariant}
Let $G$ be a compact Lie group with a bi-invariant metric. Show that the geodesics through $e$ are the one-parameter subgroups, and that $\nabla_XY = \frac12[X, Y]$ for left-invariant fields.
\end{exercise}
\begin{exercise}\label{riemannian-geometry:exercise-great-circles}
Show that the geodesics of the round sphere $S^n \subseteq \RR^{n+1}$ are the great circles $\gamma(t) = \cos(t)p + \sin(t)v$ with $|p| = |v| = 1$ and $v \perp p$ (up to affine reparametrisation).
\end{exercise}

\begin{solution}
The reflection $\rho$ of $\RR^{n+1}$ fixing the plane spanned by $p$ and $v$ pointwise preserves $S^n$ and its metric, so it is an isometry, and it maps geodesics to geodesics. The geodesic $\gamma_v$ with $\gamma_v(0) = p$, $\dot\gamma_v(0) = v$
is mapped to a geodesic with the same initial data, hence to itself by uniqueness; so $\gamma_v$ lies in the fixed set of $\rho$ on $S^n$, the great circle through $p$ and $v$. Having constant speed $1$
(Proposition~\ref{riemannian-geometry:proposition-geodesic-speed}) and starting at $p$ with velocity $v$, it is $\cos(t)p + \sin(t)v$.
\end{solution}

\begin{exercise}\label{riemannian-geometry:exercise-holonomy-product}
Show that the holonomy group of a Riemannian product $M_1 \times M_2$ at $(p_1, p_2)$ is $\mathrm{Hol}_{p_1}(M_1) \times \mathrm{Hol}_{p_2}(M_2)$, acting on $T_{p_1}M_1 \oplus T_{p_2}M_2$.
\end{exercise}

\begin{solution}
The Levi-Civita connection of the product metric is the direct sum of those of the factors, since the direct sum connection is metric and torsion-free. So a vector $(v_1, v_2)$ is transported along a curve $(c_1, c_2)$ to
$(P_{c_1}v_1, P_{c_2}v_2)$. Loops at $(p_1, p_2)$ are exactly the pairs of loops at $p_1$ and $p_2$, which gives the claim.
\end{solution}

\begin{exercise}\label{riemannian-geometry:exercise-trivial-holonomy}
Let $(M, g)$ be connected with $\mathrm{Hol}_p(g) = \{1\}$. Show that $M$ is flat and that $TM$ is trivial.
\end{exercise}

\begin{solution}
By Theorem~\ref{riemannian-geometry:theorem-holonomy-principle}, each vector of an orthonormal basis of $T_pM$ extends to a parallel vector field $E_i$; the $E_i$ are orthonormal everywhere (parallel transport is an isometry), so they
trivialise $TM$. For parallel fields $R(X, Y)E_i = \nabla_X\nabla_YE_i - \nabla_Y\nabla_XE_i - \nabla_{[X,Y]}E_i = 0$, so $R = 0$.
\end{solution}
